 % 本文件由 scripts/assemble.py 自动装配，勿手改（改 parts/ 后重跑）
% 第7章 量子霍尔态理论

% 第7章 THEORY OF QUANTUM HALL STATES
% 块界说明：本块自章首起，无 A-TAIL/JOIN。
% 书页 287（p302）末句 “...then it is easy to” 跨页延续至书页 288（p303）顶部
% “understand why it has a filling fraction ν=1/3.”，按约定半句译文留在本文件正文最后，
% 由下一块（ch7_b）的 A-TAIL 接续。
\chapter{量子霍尔态理论}

\begin{keypoints}
\item 分数量子霍尔液体代表着一些新的物质态，它们含有极为丰富的内部结构。
\item 不同的分数量子霍尔液体无法用序参量和局域算符的长程序来刻画。我们需要一套全新的理论来描述分数量子霍尔液体。
\item 分数量子霍尔液体含有无能隙的边缘激发，这些边缘激发构成所谓的手征 Luttinger 液体。边缘处的手征 Luttinger 液体也具有非常丰富的结构，这是体内丰富内部结构的一种反映。
\end{keypoints}

分数量子霍尔（FQH）效应出现在强磁场中的二维电子系统里。自 1982 年被发现以来（Tsui et al., 1982; Laughlin, 1983），对 FQH 系统的实验不断揭示出许多新现象和意外之处。所观察到的丰富分级（hierarchical）结构（Haldane, 1983; Halperin, 1984）表明，展示分数量子霍尔效应的电子态（这些态称为 FQH 液体）含有极为丰富的内部结构。事实上，FQH 液体代表着一种全新的物质态，它与前几章讨论过的对称性破缺态非常不同。要理解这种新的态，人们需要发展新的概念和新的技术。在本章中，我们将介绍 FQH 态的一般理论。对 FQH 态中新序（即拓扑序）的详细讨论将在第 8 章给出。

\section{阿哈罗诺夫--玻姆效应与分数统计}

\subsection{阿哈罗诺夫--玻姆效应——不接触而使粒子偏转}

\begin{keypoints}
\item 可缩回路（contractible loop）的局域相位与不可缩回路（non-contractible loop）的整体相位。
\item 非零的局域相位代表一种力，而非零的整体相位不对应任何经典的力。
\end{keypoints}

\begin{figure}[H]
\centering
\includegraphics[width=0.72\textwidth]{fig_7.1.pdf}
\caption*{图 7.1\quad 挖去了 $\pmb{r}=0$ 点的平面，在 $\pmb{r}=0$ 处有一根磁通管。回路 $A$ 不可缩，缠绕数 $n_w=1$；回路 $B$ 可缩。}
\end{figure}

为了给以后关于量子霍尔（QH）态中量子干涉的讨论做准备，让我们先简短地讨论一下阿哈罗诺夫--玻姆（A--B）效应。首先，我们想引入两个概念，即局域相位（local phase）与整体相位（global phase）。局域相位与可缩回路相联系。把一个带电粒子沿一条可缩回路缓慢地移动一周，会产生如下的局域相位：
\begin{equation*}
\text{局域相位} = \exp\left(\mathrm{i}e\oint\pmb{A}\,\mathrm{d}\pmb{x}\right) = \exp\left(\mathrm{i}e\int\pmb{B}\cdot\mathrm{d}\pmb{S}\right).
\end{equation*}
磁场 $\pmb{B}=0$ 当且仅当所有可缩回路的局域相位都为零。因此，不受磁力作用的带电粒子没有局域相位。

整体相位则与不可缩回路相关。为了在一个简单的背景中理解整体相位，让我们考虑一个在二维平面上运动的带电粒子。我们挖去 $\pmb{r}=0$ 点，这改变了空间的拓扑。然后在 $\pmb{r}=0$ 处放置一根无穷细的磁通管，其磁通为 $\Phi$（见图 7.1）。当粒子绕 $\pmb{r}=0$ 一周时，会产生一个相位 $\mathrm{e}^{\mathrm{i}e\Phi}$。这样的相位称为整体相位。更一般地，
\begin{gather*}
\text{整体相位} = \mathrm{e}^{\mathrm{i}e\Phi n_w}\\
n_w = \text{缠绕数}.
\end{gather*}

显然，尽管整体相位非零，粒子却感受不到任何磁力，因为在远离 $\pmb{r}=0$ 处没有磁场。我们看到，一个自由粒子（即不受力的粒子）可以具有非零的整体相位。

只有当不可缩回路存在时（即空间非单连通时），整体相位才会存在。非单连通空间的另一个例子是环面（见图 7.2）。环面上的自由粒子同样可以具有整体相位。

上述自由粒子（在 $\pmb{r}=0$ 处有一根磁通管）可以用如下哈密顿量描述：
\begin{equation*}
H = -\frac{1}{2m}\left(\pmb{\partial}-\mathrm{i}\pmb{a}\right)^2, \qquad (a_x, a_y) = \frac{1}{r^2}(-y, -x)\frac{\phi}{2\pi}
\end{equation*}

\begin{figure}[H]
\centering
\includegraphics[width=0.72\textwidth]{fig_7.2.pdf}
\caption*{图 7.2\quad 带有两个不可缩回路 $A$ 和 $B$ 的环面。}
\end{figure}

\begin{figure}[H]
\centering
\includegraphics[width=0.72\textwidth]{fig_7.3.pdf}
\caption*{图 7.3\quad 一束粒子穿过一排管子，每根管子中都有磁通。}
\end{figure}

场强为零，即
\begin{equation*}
\pmb{b} = \partial_x a_y - \partial_y a_x = 0 \ \Longrightarrow\ \text{局域相位} = 0
\end{equation*}
对于绕 $\pmb{x}=0$ 一周的回路，我们有
\begin{equation*}
\oint\pmb{a}\cdot\mathrm{d}\pmb{x} = \phi \ \Longrightarrow\ \text{整体相位} \neq 0\ .
\end{equation*}

非零的局域相位会影响经典的运动方程。自旋的贝里相就是局域相位的一个例子（见 2.3 节）。与之相反，整体相位不会影响经典的运动方程，但它会影响粒子的量子性质。

\subsection*{问题 7.1.1}

\textbf{不接触而使其偏转}\quad 尽管整体相位不产生任何经典的力，它仍然能使一个运动的粒子偏转。考虑图 7.3 中的装置：一束电荷为 $e$、动量为 $\pmb{k}$ 的粒子穿过一排不可穿透的管子。若每根管子中都有磁通 $\Phi$，试计算偏转角 $\theta$。

\begin{figure}[H]
\centering
\includegraphics[width=0.72\textwidth]{fig_7.4.pdf}
\caption*{图 7.4\quad $V_-$ 空间：挖去 $\pmb{r}=0$ 点，并把 $\pmb{r}$ 与 $-\pmb{r}$ 认同为同一个点。}
\end{figure}

\subsection{具有硬芯条件的粒子与分数统计}

\begin{keypoints}
\item 具有硬芯条件的粒子，其位形空间不是单连通的。位形空间的非平凡拓扑允许非平凡的统计。
\item 分数统计存在，而且只存在于二维空间之中。
\end{keypoints}

A--B 效应的一个具体体现是二维中的分数统计（Leinaas and Myrheim, 1977; Wilczek, 1982）。考虑二维中两个自由的、具有硬芯的全同粒子。\footnote{这里，所谓自由粒子，是指彼此分开时不受任何力作用的粒子；但这些粒子可以受到硬芯条件的约束。}其位形空间由下式给出：
\begin{equation*}
\begin{aligned}
\text{位形空间} &= \left\{(\pmb{r}_1, \pmb{r}_2)\big|_{\pmb{r}_1\neq\pmb{r}_2,\ (\pmb{r}_1,\pmb{r}_2)\sim(\pmb{r}_2,\pmb{r}_1)}\right\}\\
&= \left\{(\pmb{r}_+, \pmb{r}_-)\big|_{\pmb{r}_-\neq 0,\ (\pmb{r}_+,\pmb{r}_-)\sim(\pmb{r}_-,-\pmb{r}_-)}\right\} = V_+\otimes V_-.
\end{aligned}
\end{equation*}
其中 $\pmb{r}_+=\frac{1}{2}(\pmb{r}_1+\pmb{r}_2)$，$\pmb{r}_-=\pmb{r}_1-\pmb{r}_2$，$V_+=\{\pmb{r}_+\}$，而 $V_-=\{\pmb{r}_-|_{\pmb{r}_-\neq 0,\ \pmb{r}_-\sim-\pmb{r}_-}\}$（见图 7.4）。这里 $V_+$ 是通常的二维平面，而 $V_-$ 只是二维平面的一半，因为 $\pmb{r}_-$ 与 $-\pmb{r}_-$ 被视为同一个点。同时，$\pmb{r}_-=0$ 点已从 $V_-$ 中挖去。

如上一节所指出的，“自由”意味着局域相位为零。由于两粒子位形空间不是单连通的（$V_-$ 非单连通），整体相位存在，并且我们有权自由选择整体相位。注意，$V_-$ 中的不可缩回路由缠绕数（winding number）$n_w$ 标记（图 7.5）。于是，我们可以给不可缩回路指派如下整体相位：
\begin{equation*}
\text{整体相位} = \mathrm{e}^{\mathrm{i}\theta n_w}
\end{equation*}

现在很清楚，在量子力学中存在不同种类的自由全同粒子，由一个参数 $\theta$ 标记。

\begin{figure}[H]
\centering
\includegraphics[width=0.72\textwidth]{fig_7.5.pdf}
\caption*{图 7.5\quad $V_-$ 空间中缠绕数 $n_w=1$ 与 $n_w=2$ 的两个回路。}
\end{figure}

\begin{itemize}
\item 参数 $\theta$ 描述多粒子位形空间中的磁通（即整体相位）。
\item 参数 $\theta$ 决定统计性质。
\end{itemize}

注意，$n_w=1$ 的回路把 $\pmb{r}_-$ 连接到 $-\pmb{r}_-$，即把 $(\pmb{r}_1,\pmb{r}_2)$ 连接到 $(\pmb{r}_2,\pmb{r}_1)$，这对应于交换两个粒子。于是，粒子的统计如下：
\begin{equation*}
\left\{\begin{array}{ccc}
\theta = 0 & \rightarrow & \text{玻色子}\\
\theta = \pi & \rightarrow & \text{费米子}\\
\theta = \text{其他} & \rightarrow & \text{任意子}
\end{array}\right.
\end{equation*}

上述两个自由全同粒子（即两个具有统计 $\theta$ 的任意子（anyon））的哈密顿量为
\begin{gather*}
H = -\frac{1}{2m}(\partial_{\pmb{r}_+})^2 - \frac{1}{2}\frac{1}{m}(\partial_{\pmb{r}_-}-\mathrm{i}\pmb{a})^2\\
(a_x, a_y) = \frac{\theta}{\pi r^2}(-y, x) \tag{7.1.1}\\
\text{边界条件：}\ \ \phi(\pmb{r}_+,\pmb{r}_-) = \phi(\pmb{r}_+,-\pmb{r}_-)
\end{gather*}

当 $\theta=0$ 且 $\pmb{a}=0$ 时，$H$ 显然描述两个玻色子。对于 $\theta=\pi$，我们可以做一个规范变换
\begin{equation*}
\phi(\pmb{r}_+,\pmb{r}_-) \rightarrow \mathrm{e}^{\mathrm{i}f(\pmb{r}_-)}\phi(\pmb{r}_+,\pmb{r}_-), \qquad f(\pmb{r}) \equiv \arctan\left(\frac{x}{y}\right)
\end{equation*}

\begin{figure}[H]
\centering
\includegraphics[width=0.72\textwidth]{fig_7.6.pdf}
\caption*{图 7.6\quad 二维谐振子势阱中两个任意子的能级 $E$ 与简并度 $D$（相对运动的）。}
\end{figure}

由于 $\nabla f(\pmb{r}) = \frac{1}{r^2}(-y, x)$，我们发现 $\mathrm{e}^{-\mathrm{i}f}(\partial_{\pmb{r}_-}-\mathrm{i}\pmb{a})^2\mathrm{e}^{\mathrm{i}f} = (\partial_{\pmb{r}_-})^2$，于是哈密顿量变为
\begin{equation*}
\begin{aligned}
H &= -\frac{1}{2m}(\partial_{\pmb{r}_+})^2 - \frac{1}{2}\frac{1}{m}(\partial_{\pmb{r}_-})^2\\
\phi(\pmb{r}_+,\pmb{r}_-) &= -\phi(\pmb{r}_+,-\pmb{r}_-)\\
\text{或}\ \ \phi(\pmb{r}_1,\pmb{r}_2) &= -\phi(\pmb{r}_2,\pmb{r}_1)
\end{aligned}
\end{equation*}
（译注：原书第二行右端误印为 $-\phi(\pmb{r}_+,\pmb{r}_-)$，此处按文意改正。）
它描述两个自由费米子。一般地说，即使是 $N$ 个费米子，我们也可以自洽地把 $\pmb{a}$ 规范掉，并通过“边界条件”（即反对称条件）来描述费米统计。然而，对于 $N$ 任意子系统，我们无法把 $\pmb{a}$ 规范掉、再用“边界条件”来表示分数统计。

我们可以求解二维谐振子势阱中的两个任意子系统，得到如图 7.6 所示的能谱。可以清楚地看到，能谱如何从玻色情形连续地变到费米情形。

一个有趣的问题是：三维中存在任意子吗？三维中两个全同自由粒子的位形空间为 $V_+^{(3D)}\otimes V_-^{(3D)}$，其中
\begin{equation*}
V_-^{(3D)} = \left\{\pmb{r}_-\big|_{\pmb{r}_-\neq 0,\ \pmb{r}_-\sim-\pmb{r}_-}\right\}
\end{equation*}
这里 $V_-^{(3D)}$ 不是单连通的，整体相位存在。然而，$V_-^{(3D)}$ 中的不可缩回路由 $Z_2$ 标记。

\begin{figure}[H]
\centering
\includegraphics[width=0.72\textwidth]{fig_7.7.pdf}
\caption*{图 7.7\quad 两个交换回路 $C_1$ 与 $C_{-1}$ 在三维中可以通过绕连接两粒子的轴转动而连续地相互变换。}
\end{figure}

\begin{figure}[H]
\centering
\includegraphics[width=0.72\textwidth]{fig_7.8.pdf}
\caption*{图 7.8\quad 把一个通量--电荷束缚态绕另一个通量--电荷束缚态移动一周所诱导的相位。}
\end{figure}

这是因为沿一个方向的交换（由回路 $C_1$ 描述）与沿相反方向的交换（由回路 $C_{-1}$ 描述）属于同一类，因为两条回路 $C_{\pm 1}$ 可以连续地相互变形（见图 7.7）。如果我们给回路 $C_1$ 指派整体相位 $\mathrm{e}^{\mathrm{i}\theta}$，那么回路 $C_{-1}$ 将具有整体相位 $\mathrm{e}^{-\mathrm{i}\theta}$。由于两条回路 $C_{\pm 1}$ 仅相差一个局域相位为零的可缩回路，我们还有 $\mathrm{e}^{\mathrm{i}\theta}=\mathrm{e}^{-\mathrm{i}\theta}$。因此，$\theta$ 只能取如下两个可能的值：
\begin{equation*}
\theta = \left\{\begin{array}{cl}
0, & \text{玻色子}\\
\pi, & \text{费米子}
\end{array}\right.
\end{equation*}

以上的讨论表明，任意子是量子力学中的一种数学可能性。下一个问题是：任意子怎样才能出现在一些原本不含任意子的物理模型中？事实上，任意子可以作为玻色子/费米子系统中的激发出现。下面我们将考虑一个任意子作为电荷与磁通的束缚态而出现的模型。考虑一个由电荷为 1 的玻色子 $\phi$ 组成的系统。假设玻色子形成 $q$ 玻色子束缚态 $\Phi_q=(\phi)^q$，且 $\Phi_q$ 发生玻色凝聚。有效拉格朗日量为
\begin{align*}
\mathcal{L}_{\text{eff}} ={}& \frac{1}{2}\left|(\partial_\mu-\mathrm{i}q a_\mu)\Phi_q\right|^2 - \frac{1}{2}a\left|\Phi_q\right|^2 - \lambda\left|\Phi_q\right|^4 + \frac{1}{2g}\left(f_{\mu\nu}\right)^2 \tag{7.1.2}\\
&\qquad\qquad \langle\Phi_q\rangle \neq 0
\end{align*}
其中 $a_\mu$ 是与电荷玻色子耦合的 $U(1)$ 规范场。在凝聚态中，涡旋携带磁通 $\frac{2\pi}{q}$（见问题 7.1.2）。考虑涡旋 $\frac{2\pi}{q}$ 与原来的玻色子 $\phi$ 的束缚态。把一个束缚态绕另一个束缚态半周，我们得到一个相位
\begin{equation*}
\theta = \frac{1}{2}\frac{2\pi}{q} + \frac{1}{2}\frac{2\pi}{q} = \frac{2\pi}{q}
\end{equation*}
第一项来自电荷绕磁通管一周，第二项来自磁通管绕电荷一周（见图 7.8）。我们看到，$q=1$ 时束缚态是玻色子，$q=2$ 时是费米子，$q>2$ 时是任意子。

对于 BCS 超导体，$q=2$。但 $\phi$ 是费米子，因此束缚态是玻色子。目前所知唯一支持任意子激发的凝聚态系统是 FQH 系统。

\subsection*{问题 7.1.2}

在极坐标 $(r,\phi)$ 中，式 (7.1.2) 中的涡旋由
\begin{equation*}
\Phi_q = f(r)\mathrm{e}^{\mathrm{i}\phi}
\end{equation*}
描述，其中 $f(\infty)=\langle\Phi_q\rangle$。
\begin{enumerate}
\item[(a)] 证明：若 $a_\mu=0$，则涡旋的能量按 $\ln(L)$ 发散，其中 $L$ 是系统的尺寸。
\item[(b)] 求出使涡旋具有有限能量的 $\pmb{a}(r,\phi)$。
\item[(c)] 证明上述 $\pmb{a}(r,\phi)$ 的总磁通为 $2\pi/q$。
\end{enumerate}

\section{量子霍尔效应}

\subsection{整数量子霍尔效应}

\begin{keypoints}
\item 整数量子霍尔（IQH）态由填满的朗道能级描述。
\item $\nu=1$ IQH 态的多体波函数。
\item IQH 波函数的密度分布。
\end{keypoints}

让我们先讨论经典霍尔效应。考虑一个电荷 $q_e=-e$ 的电子二维气体，它在面内电场和垂直于该面的磁场的作用下运动。为了维持静态电流，电场力必须与磁场的洛伦兹力相平衡，即
\begin{equation*}
q_e\pmb{E} = q_e\pmb{v}\times\pmb{B}
\end{equation*}
在本章中，我们将选取使光速 $c=1$ 的单位。由于电流 $\pmb{j}$ 由 $\pmb{v}n$ 给出（$n$ 为密度），$\pmb{E}$ 与 $\pmb{j}$ 通过下式相联系：
\begin{equation*}
\pmb{E}\perp\pmb{j}, \qquad |\pmb{E}| = |\pmb{j}|\frac{B}{nq_ec} = |\pmb{j}|\frac{h}{q_e^2}\frac{B/(hc/q_e)}{n} = |\pmb{j}|\left(\frac{h}{q_e^2}\right)\frac{1}{\nu}
\end{equation*}
其中
\begin{equation*}
\nu \equiv \frac{nhc}{q_eB} = \frac{\text{粒子数}}{\text{磁通量子数}}
\end{equation*}
\begin{figure}[H]
\centering
\includegraphics[width=0.72\textwidth]{fig_7.9.pdf}
\caption*{图 7.9\quad 霍尔电阻 $\rho_{xy}$ 随 $B$ 的变化。}
\end{figure}
为填充因子。霍尔电阻 $\rho_{xy}=\left(\frac{h}{q_e^2}\right)\frac{1}{\nu}$ 可以在实验中直接测量。按照上述经典理论，若 $n$ 固定，则 $\rho_{xy}\propto B$。

实验上确实发现，对于二维电子气体，弱场下 $\rho_{xy}\propto B$（见图 7.9）。20 世纪 80 年代初，物理学家把二维电子气体置于强磁场（$\sim$10 特斯拉）中并冷却到极低的温度（$\sim$1 K），发现对于强 $B$，$\rho_{xy}$ 呈现平台结构，如图 7.9 所示（von Klitzing et al., 1980; Tsui et al., 1982）。这一现象称为 QH 效应。物理学家很快发现了 QH 态的许多惊人性质。最令人惊讶的是，平台恰好出现在与下式对应的那些 $\rho_{xy}$ 处：
\begin{equation*}
\nu = \underbrace{1, 2, 3, \cdots}_{\text{IQH}},\ \underbrace{\frac{1}{3}, \frac{2}{3}, \frac{2}{5}, \frac{3}{7}, \cdots}_{\text{FQH}}
\end{equation*}
$\rho_{xy}$ 的值如此精确和稳定，以致它已成为电阻的标准。

平台背后的物理稍微复杂一些，因为它涉及杂质。简而言之，如果二维电子气体在这些 $\nu$ 值处形成不可压缩液体\footnote{即所有带电激发都具有有限的能隙。}，我们就能理解这些填充因子处的平台。下面我们将看到，由于朗道能级结构，电子在 $\nu=1,2,3,\cdots$ 处确实形成不可压缩态。这使我们能够理解整数 $\nu$ 处的平台以及整数量子霍尔效应。

让我们忽略库仑相互作用，考虑强磁场 $B$ 中的二维自由电子气体。我们将忽略电子自旋，或者假设电子自旋已被极化。均匀磁场中的单个电子由下式描述：
\begin{gather*}
H = -\frac{1}{2m}(\partial_i-\mathrm{i}q_eA_i)^2, \qquad \hbar = c = 1 \tag{7.2.1}\\
(A_x, A_y) = \frac{B}{2}(-y, x), \qquad \text{（对称规范）}\\
B = \partial_xA_y - \partial_yA_x
\end{gather*}

上述哈密顿量有一个有趣的对称性质。尽管磁场是均匀的，它在直接平移 $T_{\pmb{a}}=\mathrm{e}^{-\pmb{a}\cdot\pmb{\partial}_{\pmb{x}}}$ 下却并非不变，即 $T_{\pmb{a}}HT_{\pmb{a}}^{\dagger}\neq H$。然而，在平移 $T_{\pmb{a}}$ 之后接着做规范变换
\begin{equation*}
G_{\pmb{a}} = \mathrm{e}^{\mathrm{i}\frac{1}{2}Bq_e(xa^y-ya^x)} \tag{7.2.2}
\end{equation*}
它就是不变的；或者更精确地说，
\begin{equation*}
G_{\pmb{a}}T_{\pmb{a}}H(G_{\pmb{a}}T_{\pmb{a}})^{\dagger} = H
\end{equation*}
组合变换 $G_{\pmb{a}}T_{\pmb{a}}$ 是 $H$ 的一个对称性，称为磁平移。不同方向的磁平移互不对易：
\begin{equation*}
(G_{\pmb{a}}T_{\pmb{a}})(G_{\pmb{b}}T_{\pmb{b}}) = \mathrm{e}^{-\mathrm{i}q_eB(a^xb^y-a^yb^x)}(G_{\pmb{b}}T_{\pmb{b}})(G_{\pmb{a}}T_{\pmb{a}}) \tag{7.2.3}
\end{equation*}
因此，尽管哈密顿量在 $x$、$y$ 两个方向都有磁平移对称性，$H$ 的本征态却不能用动量矢量 $(k_x,k_y)$ 来标记。注意，相位 $q_eB(a^xb^y-a^yb^x)$ 恰好是把电荷 $q_e$ 沿由 $\pmb{a}$、$\pmb{b}$ 张成的平行四边形 $P_{\pmb{a}\pmb{b}}$ 绕一周所得到的相位，即 $\mathrm{e}^{\mathrm{i}q_e\oint_{P_{\pmb{ab}}}\mathrm{d}\pmb{x}\cdot\pmb{A}}$。

为了求出式 (7.2.1) 的能级，我们引入复坐标 $z=x+\mathrm{i}y$，于是
\begin{equation*}
\left\{\begin{array}{ll}
\partial_z = \frac{1}{2}(\partial_x-\mathrm{i}\partial_y), & x = \frac{1}{2}(z+\bar{z}),\\
\partial_{\bar{z}} = \frac{1}{2}(\partial_x+\mathrm{i}\partial_y), & y = \frac{1}{2\mathrm{i}}(z-\bar{z}).
\end{array}\right.
\end{equation*}
用 $z$ 与 $\bar{z}$ 表示，我们有
\begin{equation*}
H = -\frac{1}{m}(D_zD_{\bar{z}} + D_{\bar{z}}D_z)
\end{equation*}
其中
\begin{equation*}
\left\{\begin{array}{ll}
D_z = \partial_z-\mathrm{i}q_eA_z, & A_z = \frac{1}{2}(A_x-\mathrm{i}A_y) = \frac{B}{4\mathrm{i}}\bar{z}\\
D_{\bar{z}} = \partial_{\bar{z}}-\mathrm{i}q_eA_{\bar{z}}, & A_{\bar{z}} = \frac{1}{2}(A_x+\mathrm{i}A_y) = -\frac{B}{4\mathrm{i}}z
\end{array}\right.
\end{equation*}
哈密顿量 $H$ 可以通过一个非幺正变换 $\mathrm{e}^{+|z|^2/4l_B^2}$ 来简化，其中 $l_B$ 是磁长度，定义为
\begin{equation*}
l_B^2 = \left|\frac{c\hbar}{q_eB}\right|
\end{equation*}
（注意 $2\pi l_B^2B = \frac{hc}{q_e} = \Phi_0$ 给出一个单位磁通量子。）我们发现
\begin{equation*}
H = \mathrm{e}^{-|z|^2/4l_B^2}\left[-\frac{1}{m}\left(\tilde{D}_z\tilde{D}_{\bar{z}} + \tilde{D}_{\bar{z}}\tilde{D}_z\right)\right]\mathrm{e}^{+|z|^2/4l_B^2}
\end{equation*}
其中
\begin{align*}
\tilde{D}_z &= \mathrm{e}^{+|z|^2/4l_B^2}D_z\mathrm{e}^{-|z|^2/4l_B^2} = \partial_z - \frac{q_eB}{4c}\bar{z} - \frac{1}{4l_B^2}\bar{z} = \partial_z - \frac{1}{2l_B^2}\bar{z}\\
\tilde{D}_{\bar{z}} &= \partial_{\bar{z}} + \frac{q_eB}{4c}z - \frac{1}{4l_B^2}z = \partial_{\bar{z}}
\end{align*}
这里我们假定了 $q_eB>0$。我们看到，非幺正变换 $\mathrm{e}^{+|z|^2/4l_B^2}$ 可以看作一个使 $A_{\bar{z}}=0$ 的“规范变换”。我们可以进一步把 $H$ 简化为
\begin{equation*}
H = \mathrm{e}^{-|z|^2/4l_B^2}\left[-\frac{2}{m}\left(\partial_z-\frac{1}{2l_B^2}\bar{z}\right)\partial_{\bar{z}}\right]\mathrm{e}^{+|z|^2/4l_B^2} + \frac{1}{2}\hbar\omega_c
\end{equation*}
其中 $\omega_c=q_eB/mc$ 是回旋频率。

现在我们可以容易地写出均匀磁场中电子的本征态：
\begin{equation*}
\begin{array}{ll}
\Psi_0 = \mathrm{e}^{-|z|^2/4l_B^2}f(z) & E = \frac{1}{2}\hbar\omega_c\\[5pt]
\Psi_1 = \mathrm{e}^{-|z|^2/4l_B^2}\left(\partial_z-\frac{1}{2l_B^2}\bar{z}\right)f(z) & E = \left(\frac{1}{2}+1\right)\hbar\omega_c\\[5pt]
\Psi_n = \mathrm{e}^{-|z|^2/4l_B^2}\left(\tilde{D}_z\right)^n f(z) & E = \left(\frac{1}{2}+n\right)\hbar\omega_c
\end{array}
\end{equation*}
这给出了朗道能级结构。

在上面的讨论中我们假定了 $q_eB>0$。第零朗道能级中的电子波函数是解析函数 $f(z)$ 乘以 $\mathrm{e}^{-|z|^2/4l_B^2}$。若 $q_eB<0$，则非幺正变换 $\mathrm{e}^{+|z|^2/4l_B^2}$ 将把 $A^z$ 变为零，第零朗道能级中的电子波函数将是反解析函数 $f(z^*)$ 乘以 $\mathrm{e}^{-|z|^2/4l_B^2}$。

第零朗道能级中的波函数 $\Psi_0$ 可以用基态
\begin{equation*}
\psi_m = \sqrt{N_m}\,z^m\,\mathrm{e}^{-|z|^2/4l_B^2}
\end{equation*}
展开，它们携带角动量 $m$。波函数 $\psi_m$ 呈环状。$|\psi_m| = \mathrm{e}^{m\ln|z|-|z|^2/4l_B^2}$ 的极大值位置在
\begin{equation*}
|z| = r_m = \sqrt{2m}\,l_B
\end{equation*}
我们发现第 $m$ 个态的环所围的面积为 $\pi r_m^2 = 2\pi l_B^2 m$，其中含有 $m$ 个磁通量子（见图 7.10）。

\begin{figure}[H]
\centering
\includegraphics[width=0.72\textwidth]{fig_7.10.pdf}
\caption*{图 7.10\quad 第零朗道能级中的圆形轨道。}
\end{figure}

\begin{figure}[H]
\centering
\includegraphics[width=0.72\textwidth]{fig_7.11.pdf}
\caption*{图 7.11\quad $\nu=1$ 液滴的密度分布，其中前 $m$ 个能级（由粗线表示）已被填满。}
\end{figure}

因此，每一个磁通量子对应一个态，而第零朗道能级中的态数等于磁通量子数。于是，当 $\nu=1$ 时，第零朗道能级中的每个态都被一个电子占据。这给出有限的激发能隙 $\Delta=\hbar\omega_c$，并解释了实验中观察到的 $\nu=1$ 平台。不可压缩性来源于泡利不相容原理（费米统计）。

$\nu=1$ 态的多体波函数具有如下形式：
\begin{equation*}
\begin{aligned}
\Psi(z_1,\cdots z_N) &= \mathcal{A}\left[(z_1)^0(z_2)^1(z_3)^2\cdots\right]\mathrm{e}^{-\sum|z_i|^2/4l_B^2}\\
&= \det\left|\begin{array}{ccc}
1 & 1 & \ldots\\
z_1 & z_2 & \ldots\\
(z_1)^2 & (z_2)^2 & \ldots\\
\vdots & \vdots & \vdots
\end{array}\right|\mathrm{e}^{-\sum|z_i|^2/4l_B^2} = \prod_{i<j}(z_i-z_j)\,\mathrm{e}^{-\sum|z_i|^2/4l_B^2}
\end{aligned}
\end{equation*}
其中 $\mathcal{A}$ 是反对称化算符。上述波函数描述一个具有均匀密度的圆形液滴（图 7.11），因为每个电子占据相同的面积 $\pi r_{m+1}^2-\pi r_m^2 = 2\pi l_B^2$。这个性质很重要。一个一般的 $N$ 变量波函数在 $N\to\infty$ 极限下未必具有均匀的密度；那样的波函数不具有合理的热力学极限。

\begin{figure}[H]
\centering
\includegraphics[width=0.72\textwidth]{fig_7.12.pdf}
\caption*{图 7.12\quad $\nu=1/3$ 液滴的密度分布。}
\end{figure}

\subsection*{问题 7.2.1}

求电子在势场 $V=\frac{1}{2}m\omega_0^2r^2$ 与均匀磁场 $B$ 中的能级。

\subsection*{问题 7.2.2}

\begin{enumerate}
\item[(a)] 证明：若规范变换 $G_{\pmb{a}}$ 由式 (7.2.2) 给出，则磁平移 $G_{\pmb{a}}T_{\pmb{a}}$ 是均匀磁场中哈密顿量 (7.2.1) 的一个对称性。
\item[(b)] 证明磁平移满足代数关系 (7.2.3)。
\end{enumerate}

\subsection{分数量子霍尔效应}

\begin{keypoints}
\item FQH 态的 Laughlin 波函数。
\item 等离子体类比与 Laughlin 态的密度分布。
\end{keypoints}

当 $0<\nu<1$ 时，第零朗道能级只是部分填充。对自由电子而言，这给出一个巨大的基态简并度 ${\sim}\,\frac{N_\phi!}{N!(N_\phi-N)!}$（其中 $N_\phi$ 是第零朗道能级中的态数，它同时也是磁通量子数；$N$ 是电子数）。因此，实验上观察到的 $\nu=1/3$ 态必定来自相互作用。这里有两个问题。第一，相互作用怎样才能产生能隙（不可压缩性）？第二，为什么 $\nu=1/3$ 是特别的？

Laughlin 用如下尝试波函数回答了上述两个问题：
\begin{equation*}
\Psi_3 = \prod_{i<j}(z_i-z_j)^3\,\mathrm{e}^{-\sum|z_i|^2/4l_B^2}
\end{equation*}
波函数 $\Psi_3$ 之所以好，有如下几条理由：
\begin{enumerate}
\item[a)] $\Psi_3$ 对每一对电子都有一个三阶零点 $(z_i-z_j)^3$。这对排斥型相互作用有利。
\item[b)] $\Psi_3$ 描述一个密度均匀的圆形液滴（图 7.12），其填充为 $\nu=1/3$。
\item[c)] $\Psi_3$ 是一个不可压缩态。
\end{enumerate}

由于 $|\Psi_3|$ 具有旋转不变性，显然 $|\Psi_3|$ 描述一个圆形液滴。如果我们再假设液滴还具有均匀的密度，那么就很容易
% ——本块到此为止：句子未完，接下一块 ch7_b（书页 288 / p303 顶部 “understand why it has a filling fraction ν=1/3.”），
% 由 ch7_b 的 A-TAIL 补出“理解它为何具有填充因子 ν=1/3。”
理解为什么它具有填充因子 $\nu=1/3$。这是因为 $z_i$（固定 $i$）的最高次幂是 $3(N-1)$。角动量为 $3(N-1)$ 的轨道的半径为 $r_{\max}=\sqrt{2\times 3(N-1)}\,l_B$。这里 $r_{\max}$ 也是液滴的半径，是 $\nu=1$ 液滴半径 $\sqrt{2\times(N-1)}\,l_B$ 的 $\sqrt{3}$ 倍。因此，对 $\Psi_3$ 有 $\nu=\frac{1}{3}$。

为了理解为什么 $\Psi_3$ 具有均匀的密度，让我们考虑下面电子位置的联合概率分布：
\begin{equation*}
P(z_1\cdots z_N)\propto|\Psi_m(z_1\cdots z_N)|^2=\mathrm{e}^{-\beta V(z_1\cdots z_N)}
\tag{7.2.4}
\end{equation*}
其中我们已把 $\Psi_3$ 态推广为 $\Psi_m$ 态
\begin{equation*}
\Psi_m=\prod_{i<j}(z_i-z_j)^m\,\mathrm{e}^{-\sum|z_i|^2/4l_B^2}
\end{equation*}
我们看到，为使 $\Psi_m$ 全反对称，$m$ 必须是奇整数。在式 (7.2.4) 中选择 $T=\frac{1}{\beta}=\frac{m}{2}$，我们可以把
\begin{equation*}
V=-m^2\sum_{i<j}\ln|z_i-z_j|+\frac{m}{4l_B^2}\sum_i|z_i|^2
\end{equation*}
视为由带"电荷" $m$ 的粒子组成的二维等离子体的势，而把 $P(z_1\cdots z_N)$ 视为该等离子体的概率分布。"电荷"为 $m_1$ 与 $m_2$ 的两个粒子之间的势是 $-m_1m_2\ln(r)$，力是 $m_1m_2/r$。因此，要理解 $\Psi_m$ 的密度分布，我们只需计算等离子体的"电荷"分布。

在 $V$ 中我们看到，等离子体中的每个粒子感受到一个势 $m|z|^2/4l_B^2$，它可以视为由一个背景"电荷"分布产生的势。实际上，这样的势由"电荷"密度为 $\rho_\phi=1/2\pi l_B^2$ 的均匀背景"电荷"产生，这正好就是磁通量子的密度。为看出这一点，我们注意到，这样的背景"电荷"作用在半径 $r$ 处的"电荷" $m$ 上的力，是"电荷" $m$ 与半径 $r$ 之内总"电荷"之间的力。该力由 $F=m\rho_\phi\pi r^2/r=m\rho_\phi\pi r$ 给出。这给出一个势 $\frac{m}{2}\pi\rho_\phi r^2$，正好就是上面的势。由于等离子体的完全屏蔽性质，等离子体倾向于"电荷"中性。等离子体的"电荷"密度等于背景"电荷"密度 $mn_e=\rho_\phi$，其中 $n_e$ 是电子密度。我们得到填充因子 $\nu=n_e/\rho_\phi=\frac{1}{m}$。

为什么 $\Psi_m$ 是不可压缩态？注意到 $\Psi_m$ 的总角动量（$z_i$ 的总幂次）为 $M_0=m\frac{N(N-1)}{2}$。$\Psi_m$ 是具有最小总角动量的态，它在每一对电子之间都有 $m$ 阶零点。压缩液滴会减小总角动量并产生更低阶的零点，这要耗费能量。我们预期该能隙为库仑能 $e^2/l_B$ 的量级。上述论断没有数学证明，但我们有大量的数值证据。

我们想指出，虽然对具有库仑相互作用的电子而言 $\Psi_3$ 只是近似基态，但存在一个理想哈密顿量，使 $\Psi_3$ 成为它的精确基态。理想哈密顿量具有形式
\begin{equation*}
H=\frac{1}{2m}\sum_i(\partial_{\pmb{x}_i}-\mathrm{i}q_e\pmb{A})^2-\sum_{i<j}UV(\pmb{r}_i-\pmb{r}_j),\quad V(\pmb{r})=\partial_{\pmb{r}}^2\delta(\pmb{r}).
\tag{7.2.5}
\end{equation*}
可以证明，对任意态 $\psi$ 有 $\langle\psi|H|\psi\rangle\geqslant\frac{1}{2}N\hbar\omega_c$，而对 $\Psi_3$ 有 $\langle\Psi_3|H|\Psi_3\rangle=\frac{1}{2}N\hbar\omega_c$。因此 $\Psi_3$ 确实是 $H$ 的精确基态。数值计算还表明，对理想哈密顿量而言 $\Psi_3$ 是不可压缩的。

\subsection*{问题 7.2.3}

证明：对理想哈密顿量 (7.2.5) 有 $\langle\Psi_3|H|\Psi_3\rangle=\frac{1}{2}N\hbar\omega_c$。为 $\nu=1/m$ 的 Laughlin 态 $\Psi_m$ 找一个理想哈密顿量。

\subsection*{问题 7.2.4}

双层 FQH 态由波函数
\begin{equation*}
\Psi_{lmn}=\prod_{i<j}(z_i-z_j)^l\prod_{i<j}(w_i-w_j)^m\prod_{i,j}(z_i-w_j)^n\,\mathrm{e}^{-\sum|z_i|^2/4l_B^2}\,\mathrm{e}^{-\sum|w_i|^2/4l_B^2}
\tag{7.2.6}
\end{equation*}
描述，其中 $z_i$ 是第一层中电子的坐标，$w_i$ 是第二层中电子的坐标。这样的态称为 $(lmn)$ 态。设 $N_1$ 与 $N_2$ 为两层中的电子数。求使两层中电子液滴尺寸相等的比值 $N_1/N_2$。求两层中电子的总填充因子。（提示：可先假设 $n=0$ 以简化问题。）

\subsection{具有分数电荷与分数统计的准粒子}

\begin{keypoints}
\item 由等离子体类比得到的分数电荷与分数统计。
\end{keypoints}

不可压缩基态 $\Psi_m$ 之上的准空穴激发由多体波函数
\begin{equation*}
\Psi^h(\xi,\xi^*)=\sqrt{C(\xi,\xi^*)}\prod_i(\xi-z_i)\Psi_m,
\end{equation*}
描述，其中 $\xi$ 是准空穴的位置，$C(\xi,\xi^*)$ 是归一化系数。计算准空穴电荷最简单的方式，是先从 $\Psi_m$ 态中移除一个电子，得到波函数 $\prod_i(\xi-z_i)^m\Psi_m$，它在 $\xi$ 处有一个电荷为 $e$ 的空穴。然而，$\prod_i(\xi-z_i)^m\Psi_m$ 也可以视为在 $\xi$ 处有 $m$ 个准空穴的波函数。因此，准空穴携带的电荷是 $e/m$。

\begin{figure}[H]
\centering
\includegraphics[width=0.72\textwidth]{fig_7.13.pdf}
\caption*{图 7.13\quad 电子的电荷与测试电荷必须中和背景电荷。}
\end{figure}

计算准空穴电荷更直接的途径是等离子体类比。准空穴态 $\Psi_h$ 中电子位置的分布为
\begin{equation*}
P^h_\xi(\{z_i\})=\bigl|\Psi^h(\xi,\xi^*)\bigr|^2=\mathrm{e}^{-\beta[V(z_i)-m\ln|\xi-z_i|]}
\end{equation*}
因此，电子感受到一个附加的背景势 $-m\ln|\xi-z_i|$。这个势由一个单位测试电荷产生。背景电荷在 $\xi$ 处包含一个额外的单位测试电荷。为维持电荷中性，$\xi$ 附近必须少 $1/m$ 个电子（见图 7.13）。这样，准空穴携带电荷 $e/m$。

为了理解准空穴的分数统计（Arovas et al., 1984），让我们从下面几个准粒子的有效拉格朗日量出发：
\begin{equation*}
Z=\int\mathcal{D}[\pmb{x}_i(t)]\;\mathrm{e}^{\mathrm{i}\int\mathrm{d}t\,\sum_i[-E_0+\frac{1}{2}m\dot{\pmb{x}}_i^2+\pmb{a}(\pmb{x}_1,\pmb{x}_2,\ldots)\cdot\dot{\pmb{x}}_i]}
\end{equation*}
有效规范势 $\pmb{a}$ 决定统计。正如 2.3 节讨论的自旋系统那样，$\pmb{a}$ 由相干态 $|\Psi^h(\xi(t),\xi^*(t))\rangle$ 的贝里相决定。

让我们先只考虑单个准粒子。对于准粒子 $\xi(t)$ 的绝热运动，$\Psi^h(\xi(t),\xi^*(t))$ 的相位变化给出贝里相。对小的 $\Delta t$，贝里相为
\begin{equation*}
\langle\Psi^h(\xi(t+\Delta t),\xi^*(t+\Delta t))|\Psi^h(\xi(t),\xi^*(t))\rangle=\mathrm{e}^{\mathrm{i}\Delta t\,\pmb{a}\cdot\dot{\pmb{x}}}.
\end{equation*}
于是 $\pmb{a}$ 由
\begin{equation*}
\mathrm{i}\langle\Psi^h(\xi(t),\xi^*(t))|\frac{\mathrm{d}}{\mathrm{d}t}|\Psi^h(\xi(t),\xi^*(t))\rangle=\pmb{a}\cdot\dot{\pmb{x}}
\end{equation*}
给出。由于 $\Psi^h$ 只通过 $\xi(t)$ 依赖 $t$，我们有
\begin{align*}
\pmb{a}\cdot\dot{\pmb{x}}&=\mathrm{i}\langle\Psi^h(\xi(t),\xi^*(t))|\frac{\mathrm{d}}{\mathrm{d}t}|\Psi^h(\xi(t),\xi^*(t))\rangle\\
&=\frac{\mathrm{d}\xi}{\mathrm{d}t}\,\mathrm{i}\langle\Psi^h(\xi,\xi^*)|\frac{\partial}{\partial\xi}|\Psi^h(\xi,\xi^*)\rangle+\frac{\mathrm{d}\xi^*}{\mathrm{d}t}\,\mathrm{i}\langle\Psi^h(\xi,\xi^*)|\frac{\partial}{\partial\xi^*}|\Psi^h(\xi,\xi^*)\rangle
\end{align*}
这样，如果我们写出 $\pmb{a}\cdot\mathrm{d}\pmb{x}=a_\xi\mathrm{d}\xi+a_{\xi^*}\mathrm{d}\xi^*$，就有\footnote{我们用到了
\[
\begin{aligned}
&\mathrm{i}\langle\Psi^h(\xi,\xi^*)|\frac{\partial}{\partial\xi}|\Psi^h(\xi,\xi^*)\rangle=\mathrm{i}\langle\Psi(\xi)|\frac{1}{\sqrt{C(\xi,\xi^*)}}\frac{\partial}{\partial\xi}\frac{1}{\sqrt{C(\xi,\xi^*)}}|\Psi(\xi)\rangle\\
&\quad=\mathrm{i}\int\prod_i\mathrm{d}^2z_i\;\Psi^*(\xi^*)\frac{1}{\sqrt{C}}\frac{\partial}{\partial\xi}\Bigl(\frac{1}{\sqrt{C}}\Psi(\xi)\Bigr)\\
&\quad=\mathrm{i}\frac{\partial}{\partial\xi}\int\prod_i\mathrm{d}^2z_i\;\Psi^*(\xi^*)\frac{1}{C}\Psi(\xi)-\mathrm{i}\int\prod_i\mathrm{d}^2z_i\;\Psi^*(\xi^*)\Bigl(\frac{\partial}{\partial\xi}\frac{1}{\sqrt{C}}\Bigr)\frac{1}{\sqrt{C}}\Psi(\xi)\\
&\quad=-\mathrm{i}\sqrt{C}\frac{\partial}{\partial\xi}\frac{1}{\sqrt{C}}
\end{aligned}
\]}
\begin{equation*}
a_\xi=\mathrm{i}\langle\Psi^h(\xi,\xi^*)|\frac{\partial}{\partial\xi}|\Psi^h(\xi,\xi^*)\rangle=-\mathrm{i}\sqrt{C}\frac{\partial}{\partial\xi}\frac{1}{\sqrt{C}}
\end{equation*}
这里 $C(\xi,\xi^*)=\langle\Psi(\xi)|\Psi(\xi)\rangle$，而 $\Psi(\xi)=\prod(\xi-z_i)\Psi_m$ 是 $\xi$ 的解析函数。上述结果的关键之点是，$\Psi$（$\xi$ 的解析函数）的归一化给出贝里相：
\begin{equation*}
a_\xi=+\frac{\mathrm{i}}{2}\frac{\partial}{\partial\xi}\ln C,\qquad a_{\xi^*}=-\frac{\mathrm{i}}{2}\frac{\partial}{\partial\xi^*}\ln C
\end{equation*}
归一化 $C(\xi,\xi^*)$ 可以用等离子体类比算出。我们注意到，在加入一个代表测试"电荷"与背景"电荷"之间相互作用的项 $\mathrm{e}^{-|\xi|^2/4ml_B^2}$ 之后，\footnote{注意 $\mathrm{e}^{-|\xi|^2/4l_B^2}$ 对应于一个电子与背景"电荷"之间的相互作用。一个电子对应于 $m$ 个测试"电荷"的束缚态。}
\begin{equation*}
\int\prod_i\mathrm{d}^2z_i\;\Bigl|\mathrm{e}^{-\frac{1}{4ml_B^2}|\xi|^2}\prod(\xi-z_i)\prod(z_i-z_j)^m\,\mathrm{e}^{-\frac{1}{4l_B^2}|z_i|^2}\Bigr|^2\propto\mathrm{e}^{-\beta V(\xi,\xi^*)}
\end{equation*}
给出的是在 $\xi$ 处带有一个测试"电荷"的等离子体的总能量。同样，由于等离子体的完全屏蔽性质，作用在插入的测试"电荷"上的总力为零；因而等离子体的能量不依赖于 $\xi$，即 $V(\xi,\xi^*)=\text{常数}$。我们看到
\begin{equation*}
C(\xi,\xi^*)\propto\mathrm{e}^{\frac{1}{2l_B^2}\frac{1}{m}|\xi|^2}
\end{equation*}
这给出
\begin{equation*}
a_\xi=\mathrm{i}\frac{1}{m}\frac{1}{4l_B^2}\xi^*,\qquad a_{\xi^*}=-\mathrm{i}\frac{1}{m}\frac{1}{4l_B^2}\xi
\end{equation*}
上述矢量势 $\pmb{a}$ 正比于磁场的矢量势：
\begin{equation*}
a_\xi=-\frac{1}{m}q_eA_z,\qquad a_{\xi^*}=-\frac{1}{m}q_eA_{z^*}
\end{equation*}
绕环路一周，准空穴获得 A--B 相 $\oint\mathrm{d}\pmb{x}\cdot\pmb{a}=\bigl(-\frac{q_e}{m}\bigr)\oint\mathrm{d}\pmb{x}\cdot\pmb{A}$，即一个电荷为 $-q_e/m=e/m$ 的粒子的相位。

对于两个准空穴，带两个测试"电荷"的等离子体的总能量 $V$ 由
\begin{align*}
&\mathrm{e}^{-\beta V(\xi_1,\xi_1^*,\xi_2,\xi_2^*)}\\
&\quad=\Bigl|\mathrm{e}^{-\frac{1}{4l_B^2}\frac{1}{m}(|\xi_1|^2+|\xi_2|^2)}\,|\xi_1-\xi_2|^{\frac{1}{m}}\prod_{a,i}(\xi_a-z_i)\prod(z_i-z_j)^m\,\mathrm{e}^{-\frac{1}{4l_B^2}\sum|z_i|^2}\Bigr|^2
\end{align*}
给出，其中额外的项 $|\xi_1-\xi_2|^{2/m}=\mathrm{e}^{-\beta(-)\ln|\xi_1-\xi_2|}$ 表示两个测试"电荷"之间的直接相互作用，每个"电荷"为 1。同样，$V(\xi_1,\xi_1^*,\xi_2,\xi_2^*)$ 不依赖于 $\xi_1$ 与 $\xi_2$。于是
\begin{equation*}
|\Psi(\xi_1,\xi_2)|^2\propto\mathrm{e}^{\frac{1}{2l_B^2}\frac{1}{m}(|\xi_1|^2+|\xi_2|^2)}\,|\xi_1-\xi_2|^{-\frac{2}{m}}
\end{equation*}
令 $\xi=\xi_1$、$\xi_2=0$，我们发现
\begin{align*}
a_\xi&=\mathrm{i}\frac{1}{m}\frac{1}{4l_B^2}\xi^*-\frac{\mathrm{i}}{2m}\frac{\partial}{\partial\xi}\ln|\xi|^2=\mathrm{i}\frac{1}{m}\frac{1}{4l_B^2}\xi^*-\frac{\mathrm{i}}{2m}\frac{1}{\xi}\\
a_{\xi^*}&=-\mathrm{i}\frac{1}{m}\frac{1}{4l_B^2}\xi+\frac{\mathrm{i}}{2m}\frac{1}{\xi^*}
\end{align*}
这给出 $(a_x,a_y)=(-y,x)\frac{1}{r^2}\frac{1}{m}+\cdots$。当我们把准空穴 $\xi_1$ 绕 $\xi_2$ 移动一周时，准空穴获得的相位为 $\frac{e}{m}B\times\text{面积}+\frac{2\pi}{m}$。第一项来自均匀磁场。附加项 $\frac{2\pi}{m}$ 来自 $(-y,x)\frac{1}{r^2}\frac{1}{m}$（译注：原文此处印作 $\frac{1}{r}$）。正是这一项赋予准空穴统计角为 $\theta=\frac{\pi}{m}$ 的分数统计（见式 (7.1.1)）。

\subsection*{问题 7.2.5}

考虑双层 $(lmn)$ 态，即式 (7.2.6)。第一层中的准空穴由波函数
\begin{equation*}
\Psi^h_1(\xi_1)=\prod_i(\xi_1-z_i)\Psi_{lmn}(\{z_i,w_j\})
\end{equation*}
描述，第二层中的准空穴由
\begin{equation*}
\Psi^h_2(\xi_2)=\prod_i(\xi_2-w_i)\Psi_{lmn}(\{z_i,w_j\})
\end{equation*}
描述。求第一层与第二层中准空穴的分数电荷与分数统计。求第一层准空穴与第二层准空穴之间的互统计。两个非全同粒子之间的互统计定义为把一个粒子绕另一个粒子一整周所产生的贝里相。

\begin{figure}[H]
\centering
\includegraphics[width=0.72\textwidth]{fig_7.14.pdf}
\caption*{图 7.14\quad 分级 $\nu=2/7$ 与 $\nu=2/5$ 态。}
\end{figure}

\subsection{分级分数量子霍尔态——Laughlin 理论的推广}

\begin{keypoints}
\item 分级 FQH 态等价于准粒子的 Laughlin 态。
\end{keypoints}

要构造填充因子不是 $1/m$ 的更一般的 FQH 态，我们可以从 $\nu=1/m$ 的 Laughlin 态出发，加入准空穴或准粒子来改变平均电子密度。当其密度达到某一数值时，准空穴/准粒子气体自身也能形成一个 Laughlin 态。所得的态就是一个分级 FQH 态（见图 7.14）。

当 $\nu=1/3$ 态之上的电荷为 $e/3$ 的准空穴凝聚成一个 Laughlin 态时，我们得到 $\nu=2/7$ 的 FQH 态。波函数具有形式
\begin{align*}
&\Psi(z_1\cdots z_N)\\
&\quad=\int\prod_i\mathrm{d}^2\xi_i\;\prod(z_i-z_j)^3\,\mathrm{e}^{-\frac{1}{4l_B^2}\sum|z_i|^2}\prod(\xi_i-z_j)(\xi_i^*-\xi_j^*)^2\,\mathrm{e}^{-\frac{1}{4l_B^2}\frac{1}{3}|\xi_i|^2}
\end{align*}
当 $\nu=1/3$ 态之上的电荷为 $e/3$ 的准粒子凝聚成一个 Laughlin 态时，我们得到 $\nu=2/5$ 的 FQH 态。波函数为
\begin{equation*}
\Psi_c(z_i-z_N)=\int\prod_i\mathrm{d}^2\xi_i\;\prod_{i<j}(\xi_i-\xi_j)^2(\xi_i^*-2\partial_{z_i})\prod(z_i-z_j)^3\,\mathrm{e}^{-\frac{1}{4l_B^2}\sum|z_i|^2}
\end{equation*}
（译注：原书此式末尾的指数印作 $-\frac{1}{4l_B^2}|z_i|$，疑脱求和号与平方。）

我们可以用等离子体类比计算上述两个态的性质，但计算很复杂。下一节我们将推导 FQH 态的有效场论，并用有效理论计算上述分级 FQH 态的性质。有效场论大大简化了计算，我们能容易地得到 FQH 态的填充因子与准粒子的量子数。

\section{分数量子霍尔液体的有效理论}

\begin{keypoints}
\item FQH 态的低能有效理论是 $U(1)$ Chern--Simons 理论，它们捕捉 FQH 态的普适性质。
\end{keypoints}

我们已经见过几种不同的 FQH 态。一方面，这些 QH 态具有不同的分数电荷与不同的分数统计，这暗示不同的 FQH 态属于不同的量子相。另一方面，这些 FQH 态又都具有相同的对称性，我们无法用破缺的对称性来区分它们。实际上，FQH 态包含一种新型的序——拓扑序（见第 8 章）。我们需要用新的工具来刻画拓扑序。

系统研究拓扑序的一种途径是为 FQH 液体构造低能有效理论。有效理论能够捕捉 FQH 液体的普适性质，并为如何刻画和标记 FQH 液体中不同的拓扑序提供线索。下面我们将介绍一种构造有效理论的方法，它与 Haldane 和 Halperin 提出的分级构造密切相关（Haldane, 1983; Halperin, 1984）。

刻画 FQH 液体内部结构的首次尝试由 Girvin 和 MacDonald (1987) 提出，他们证明 Laughlin 态在一个非局域算符中具有非对角长程序。这一观察导致用 Ginzburg--Landau--Chern--Simons 有效理论来描述 FQH 液体（Read, 1989; Zhang et al., 1989）。这些进展产生了许多有趣而重要的结果，并加深了对 FQH 液体的理解。然而，这里我将尝试从更一般的观点出发描述 FQH 液体的内部结构。看来，FQH 液体的某些内部结构（特别是在所谓非阿贝尔 FQH 液体中的那些）无法用 Ginzburg--Landau--Chern--Simons 有效理论及相关的非对角长程序来描述。因此，我们需要发展更普遍的概念与表述，如拓扑序与拓扑场论，来描述 FQH 液体中的内部结构（见第 8 章）。在我看来，非对角长程序的概念并没有抓住 FQH 液体内部结构的本质。这里我们将不使用非对角长程序来发展有效理论（Blok and Wen, 1990a,b）。所得的有效理论只含纯 Chern--Simons 项。纯 Chern--Simons 形式更为紧凑，也更清楚地揭示分级 FQH 态的内部结构。

\subsection{Laughlin 态的有效理论}

\begin{keypoints}
\item FQH 液体的流体动力学与有效 Chern--Simons 理论。
\end{keypoints}

本节中我们只考虑单层自旋极化的 QH 系统。为构造分级态的有效理论，我们先设法得到 Laughlin 态的有效理论，然后在下一节用分级构造得到分级态的有效理论。

考虑磁场中的一个电子系统。为了涵盖更一般的情形，我们假设电子可以具有玻色统计或费米统计。系统的拉格朗日量具有如下形式（第一次量子化的版本）：
\begin{equation*}
\begin{aligned}
\mathcal{L}&=-eA^iJ^i+\text{动能/势能}\\
&=eA_iJ^i+\text{动能/势能}
\end{aligned}
\tag{7.3.1}
\end{equation*}
其中
\begin{equation*}
\pmb{J}(\pmb{x})=\sum_i\pmb{v}_i\delta(\pmb{x}-\pmb{x}_i),\qquad J^0(\pmb{x})=\sum_i\delta(\pmb{x}-\pmb{x}_i)
\tag{7.3.2}
\end{equation*}
分别是电子的流与密度，$(\pmb{x}_i,\pmb{v}_i)$ 是第 $i$ 个电子的位置与速度。动能/势能由 $\sum_i\frac{1}{2}mv_i^2+\sum_{i<j}V(\pmb{x}_i-\pmb{x}_j)$ 给出，其具体形式对我们并不重要。我们还假设了每个电子的电荷为 $-e$，光速为 $c=1$。

在流体动力学方法中，我们假设低能集体模式可以用密度与流 $J^\mu$ 描述，而低能有效理论具有形式 $\mathcal{L}(A_\mu,J^\nu)=eA_iJ^i+\mathcal{L}'(J^\mu)$。从 5.1.3 节与 5.3.3 节的讨论我们看到，有时一个态的低能激发可能太多，以至于无法用单一密度模式来描述。由于有能隙的 FQH 态的低能激发比超流体还要少，可以合理地假设单一密度模式足以描述 FQH 态的低能涨落。

在填充因子 $\nu=1/m$ 处（$m$ 对玻色型电子为偶整数，对费米型电子为奇整数），电子系统的基态由 Laughlin 波函数（Laughlin, 1983）给出：
\begin{equation*}
\Bigl[\prod(z_i-z_j)^m\Bigr]\mathrm{e}^{-\frac{1}{4l_B^2}\sum|z_i|^2}
\tag{7.3.3}
\end{equation*}
（这里我们已假设 $B<0$，于是 $(-e)B>0$。）为构造有效理论，我们注意到态 (7.3.3) 是不可压缩流体，其密度与磁场绑定，即 $J^0=(-e)\nu B/2\pi$。再结合有限的霍尔电导 $\sigma_{xy}=\frac{\nu e^2}{h}=\frac{\nu e^2}{2\pi\hbar}=\frac{\nu e^2}{2\pi}$，我们发现电子数流 $J^\mu$ 对电磁场的变化有如下响应：
\begin{equation*}
-e\delta J^\mu=-\sigma_{xy}\varepsilon^{\mu\nu\lambda}\partial_\nu\delta A_\lambda=-\frac{\nu e^2}{2\pi}\varepsilon^{\mu\nu\lambda}\partial_\nu\delta A_\lambda
\tag{7.3.4}
\end{equation*}
我们这样选择有效拉格朗日量 $\mathcal{L}(A_\mu,J^\nu)$，使它产生运动方程 (7.3.4)。为方便起见，引入一个 $U(1)$ 规范场 $a_\mu$ 来描述电子数流：
\begin{equation*}
J^\mu=\frac{1}{2\pi}\partial_\nu a_\lambda\varepsilon^{\mu\nu\lambda}
\end{equation*}
这样定义的流自动满足守恒定律。于是，产生式 (7.3.4) 的有效拉格朗日量取如下形式：
\begin{equation*}
\mathcal{L}=\Bigl[-m\frac{1}{4\pi}a_\mu\partial_\nu a_\lambda\varepsilon^{\mu\nu\lambda}+\frac{e}{2\pi}A_\mu\partial_\nu a_\lambda\varepsilon^{\mu\nu\lambda}\Bigr]
\tag{7.3.5}
\end{equation*}

式 (7.3.5) 只描述基态对外电磁场的线性响应。要对拓扑流体（例如 FQH 液体）有更完整的描述，我们需要把电子激发引入有效理论。特别是，我们要确保有效理论中包含一个携带与电子相同量子数的激发。

\subsection{有效理论中的电子与准粒子激发}

\begin{keypoints}
\item 单凭有效 Chern--Simons 理论并不是对 FQH 态的完整描述。为得到完整的有效理论，我们还需要指明电子算符。
\item 电子与准粒子/准空穴之间所要求的平凡互统计决定了准粒子与准空穴的量子数。
\end{keypoints}

让我们引入一个携带 $a_\mu$ 电荷 $l$ 的粒子。在有效理论 (7.3.5) 中，这样的粒子对应于如下源项：
\begin{equation*}
la_0\delta(\pmb{x}-\pmb{x}_0)
\tag{7.3.6}
\end{equation*}
该源项将产生一个电荷为
\begin{equation*}
Q=-le/m.
\tag{7.3.7}
\end{equation*}
的激发。这可以从运动方程 $\delta\mathcal{L}/\delta a_0=0$ 看出，由此方程得到
\begin{equation*}
J_0=\frac{1}{2\pi}\varepsilon_{ij}\partial_i a_j=-\frac{e}{2\pi m}B+\frac{l}{m}\delta(\pmb{x}-\pmb{x}_0)
\end{equation*}
右端第一项表明填充因子 $\nu\equiv 2\pi\frac{J_0}{-eB}$ 确实为 $\nu=1/m$，第二项则对应于与该激发相联系的电子密度的增加。

我们还看到，由源项 (7.3.6) 产生的激发带有 $l/m$ 个单位的 $a_\mu$ 通量。这样，如果有两个分别携带 $a_\mu$ 电荷 $l_1$ 与 $l_2$ 的激发，那么把一个激发绕另一个激发移动一周将诱导一个相位 $2\pi\times(a_\mu\text{-通量量子数})\times(a_\mu\text{ 电荷})$，\footnote{细心的读者可能注意到，这两个激发既带有 $a_\mu$ 电荷又带有 $a_\mu$ 通量。因此，把一个激发绕另一个激发移动一周所诱导的相位似乎应收到两项贡献：一项来自电荷绕通量运动，另一项来自通量绕电荷运动，如图 7.8 所描述的那样。这将给出相位 $2\pi\times\frac{l_1}{m}\times l_2+2\pi\times\frac{l_2}{m}\times l_1$，是式 (7.3.8) 中结果的两倍。然而，更仔细的计算表明，图 7.8 的图像并不适用于由 Chern--Simons 项诱导的电荷--通量束缚态。式 (7.3.8) 中的朴素结果恰好是正确的（见问题 7.3.1）。}即
\begin{equation*}
2\pi\times\frac{l_1}{m}\times l_2
\tag{7.3.8}
\end{equation*}
如果 $l_1=l_2\equiv l$，那么这两个激发是全同的。把它们互换将诱导式 (7.3.8) 中相位的一半，即
\begin{equation*}
\theta=\pi\frac{l^2}{m}
\tag{7.3.9}
\end{equation*}
这里的 $\theta$ 正是携带 $l$ 个单位 $a_\mu$ 电荷的激发的统计角。

我们的电子携带电荷 $-e$。从上面的讨论我们看到，电荷为 $-e$ 的激发对应于携带 $m$ 个单位 $a_\mu$ 电荷的粒子。这样的激发具有统计角 $\theta=\pi m$（见式 (7.3.9)）；因此，$m$ 为偶数时它是玻色子，$m$ 为奇数时它是费米子。于是，对玻色型电子取偶数 $m$ 的情形，以及对费米型电子取奇数 $m$ 的情形，我们可以把携带 $m$ 个单位 $a_\mu$ 电荷的激发认定为构成 FQH 液体的电子。具有电子的电荷与统计的激发的存在，支持了我们的 Chern--Simons 有效理论的正确性。

我们想强调，在有效理论中认定基本电子是非常重要的。正是这一认定，连同有效拉格朗日量，提供了对 FQH 液体拓扑性质的完整描述。我们将在下面看到，这一认定使我们能够确定准粒子激发的分数电荷与分数统计。

位置由复数 $\xi=x+\mathrm{i}y$ 描述的准空穴激发，是通过用 $\prod_i(\xi-z_i)$ 乘以基态波函数 (7.3.3) 产生的。注意到当一个电子绕准空穴一周时，波函数的相位改变 $2\pi$。这个 $2\pi$ 相位意味着电子与准空穴具有平凡的互统计。一般地，电子与一个允许的激发必须具有平凡的互统计，这样该激发的电子波函数才是单值的。

现在让我们试着通过插入一个携带 $l$ 个单位 $a_\mu$ 电荷的源项来产生激发。把一个电子绕这样的激发移动一周将诱导相位 $2\pi l$（见式 (7.3.8)）。电子波函数的单值性要求这样的相位是 $2\pi$ 的整数倍。因此 $a_\mu$ 电荷必须量子化为整数，也只有这些电荷对应的才是允许的激发。

从式 (7.3.7) 给出的激发电荷我们发现，$l=-1$ 对应于基本准空穴激发，而 $l=1$ 对应于基本准粒子激发。准粒子激发携带电荷 $-e/m$，准空穴携带电荷 $e/m$。从式 (7.3.9) 可见，二者都具有统计 $\theta=\pi/m$。我们看到，有效理论重现了 Laughlin 态中准粒子的著名结果（Arovas et al., 1984）。包含准粒子激发的完整有效理论为
\begin{equation*}
\begin{aligned}
\mathcal{L}&=\Bigl[-m\frac{1}{4\pi}a_\mu\partial_\nu a_\lambda\varepsilon^{\mu\nu\lambda}+\frac{e}{2\pi}A_\mu\partial_\nu a_\lambda\varepsilon^{\mu\nu\lambda}\Bigr]\\
&\quad+la_\mu j^\mu+\text{动能/势能}
\end{aligned}
\tag{7.3.10}
\end{equation*}
其中 $j^\mu$ 是准粒子的流，其形式由式 (7.3.2) 给出。式 (7.3.10) 连同对 $l$ 的量子化条件，是一个完整的低能有效理论，它捕捉了 $1/m$ Laughlin 态的拓扑性质。

\subsection*{问题 7.3.1}

证明式 (7.3.8)。（提示：可以从下面带两个激发的有效理论出发：
\begin{equation*}
\mathcal{L}=\Bigl[-\frac{1}{2}\frac{m}{2\pi}a_\mu\partial_\nu a_\lambda\varepsilon^{\mu\nu\lambda}+j^\mu a_\mu\Bigr]
\end{equation*}
其中 $j^\mu=j_1^\mu+j_2^\mu$，而
\begin{equation*}
\begin{aligned}
j_1^0(\pmb{x},t)&=l_1\delta(\pmb{x}-\pmb{x}_1(t)),&\qquad j_2^0(\pmb{x},t)&=l_2\delta(\pmb{x}-\pmb{x}_2(t)),\\
j_1^i(\pmb{x},t)&=l_1\dot{x}_1^i\delta(\pmb{x}-\pmb{x}_1(t)),&\qquad j_2^i(\pmb{x},t)&=l_2\dot{x}_2^i\delta(\pmb{x}-\pmb{x}_2(t)).
\end{aligned}
\end{equation*}
这里 $j^\mu$ 是两个激发的总流，$\pmb{x}_{1,2}(t)$ 是两个激发的位置。然后可以积掉 $a_\mu$，得到
\begin{equation*}
\int\mathrm{d}^3x\;\frac{1}{2}j^\mu\frac{1}{\frac{m}{2\pi}\partial_\lambda\epsilon^{\mu\lambda\nu}}j^\nu
\end{equation*}
交叉项可以写成
\begin{equation*}
\int\mathrm{d}^3x\;j_1^\mu\frac{1}{\frac{m}{2\pi}\partial_\lambda\epsilon^{\mu\lambda\nu}}j_2^\nu=\int\mathrm{d}^3x\;j_1^\mu f_\mu
\end{equation*}
其中 $f_\mu$ 满足
\begin{equation*}
\frac{m}{2\pi}\partial_\lambda\epsilon^{\mu\lambda\nu}f_\nu=j_2^\mu
\end{equation*}
假设 $\pmb{x}_2=\dot{\pmb{x}}_2=0$，即可求出 $f_i$。）

\subsection*{问题 7.3.2}

前面我们一直关注 Laughlin 态的拓扑性质。要对 Laughlin 态的动力学性质有所了解，我们需要在有效 Chern--Simons 理论中纳入如下的 Maxwell 项：
\begin{equation*}
\mathcal{L}=-m\frac{1}{4\pi}a_\mu\partial_\nu a_\lambda\varepsilon^{\mu\nu\lambda}+\frac{1}{2g_1}e^2-\frac{1}{2g_2}b^2
\tag{7.3.11}
\end{equation*}
其中 $e$ 与 $b$ 分别是 $a_\mu$ 的电场与磁场，且 $e_i=\dot{a}_i-\partial_i a^0$，$b=\partial_1a_2-\partial_2a_1$。求由 $e$ 与 $b$ 描述的集体涨落的运动方程。求集体激发的能隙。假设 $1/m$ Laughlin 态中的能隙为 $e^2/m^2 l_B$ 的量级，准空穴的尺寸为 $l_B$ 的量级。估计 $g_1$ 与 $g_2$ 的值。

\subsection{分级分数量子霍尔态的有效理论}

\begin{keypoints}
\item 不同的分级 FQH 态可以用一个整数 $K$ 矩阵和一个电荷矢量 $\pmb{q}$ 表征。
\item FQH 态的拓扑性质，如填充因子与准粒子量子数，可以容易地由 $K$ 矩阵和电荷矢量算出。
\item 不同的 $K$ 矩阵--电荷矢量对之间的等价关系。
\end{keypoints}

为得到分级 FQH 态的有效理论，让我们从费米型电子形成的 $1/m$ Laughlin 态出发。通过产生基本准粒子（记为 $l=1$）来增大填充因子。$l=1$ 的式 (7.3.10) 描述了存在这些准粒子时的 $1/m$ 态。现在出现两幅等价的图像。

\textbf{(a)}~在平均场理论方法中，我们可以把式 (7.3.10) 中的规范场 $a_\mu$ 视为固定背景，不允许 $a_\mu$ 对插入的源项 $j^\mu$ 作出响应。在这种情形下，准粒子气体的行为像"磁"场 $b=\partial_i a_j\varepsilon_{ij}$ 中的玻色子，这从式 (7.3.10) 的第二项可以看出。这些玻色子不携带任何电荷，因为准粒子数流 $j^\mu$ 不与电磁规范势 $A_\mu$ 直接耦合。当玻色子密度满足
\begin{equation*}
j^0=\frac{1}{p_2}\frac{b}{2\pi}
\end{equation*}
其中 $p_2$ 为偶数时，玻色子具有填充因子 $\frac{1}{p_2}$。玻色子的基态又可以用一个 Laughlin 态来描述。我们得到的最终电子态正是 Haldane (1983) 所构造的第二级分级 FQH 态。

\textbf{(b)}~如果我们让 $a_\mu$ 对 $j^\mu$ 的插入作出响应，那么准粒子将被 $a_\mu$ 通量缀饰。被缀饰的准粒子携带电荷 $-e/m$，统计为 $\theta=\pi/m$。当准粒子具有密度
\begin{equation*}
j^0=\frac{1}{(p_2-\frac{\theta}{\pi})}\frac{(-e)B}{2\pi m}
\end{equation*}
其中 $p_2$ 为偶数时，准粒子将具有填充因子 $\frac{1}{(p_2-\frac{\theta}{\pi})}$。这时，准粒子系统可以形成由波函数
\begin{equation*}
\prod_{i<j}(z_i-z_j)^{p_2-\frac{\theta}{\pi}}
\end{equation*}
描述的 Laughlin 态。用这种方法得到的最终电子态同样是一个第二级分级 FQH 态。这一构造由 Halperin (1984) 首先提出。(a) 与 (b) 两种构造给出同一个分级态，它们是等价的。

下面，我们将遵循 Haldane 的分级构造来推导分级 FQH 态的 Chern--Simons 有效理论（Blok and Wen, 1990a,b; Read, 1990; Fr\"{o}hlich and Kerler, 1991; Wen and Zee, 1992a）。注意到，在假设 (a) 之下，玻色子拉格朗日量（$l=1$ 时式 (7.3.10) 的第二项）正是把外电磁场 $eA_\mu$ 换成 $a_\mu$ 后的式 (7.3.1)。于是，我们可以重复从式 (7.3.3) 到式 (7.3.10) 的步骤，来构造玻色子 Laughlin 态的有效理论。引入一个新的 $U(1)$ 规范场 $\tilde{a}_\mu$ 来描述玻色子流，我们发现玻色子有效理论具有形式
\begin{equation*}
\mathcal{L}=-\frac{p_2}{4\pi}\tilde{a}_\mu\partial_\nu\tilde{a}_\lambda\varepsilon^{\mu\nu\lambda}+\frac{1}{2\pi}a_\mu\partial_\nu\tilde{a}_\lambda\varepsilon^{\mu\nu\lambda}
\tag{7.3.12}
\end{equation*}
在式 (7.3.12) 中，新的规范场 $\tilde{a}_\mu$ 描述玻色子的密度 $j^0$ 与流 $j^i$，它由
\begin{equation*}
j^\mu=\frac{1}{2\pi}\partial_\nu\tilde{a}_\lambda\varepsilon^{\mu\nu\lambda}
\end{equation*}
给出。这把 $a_\mu$ 与玻色子流 $a_\mu j^\mu$ 的耦合约化为 $a_\mu$ 与 $\tilde{a}_\mu$ 之间的一个 Chern--Simons 项（它成为式 (7.3.12) 中的第二项）。总的有效理论（包括原来的电子凝聚体）具有形式
\begin{equation*}
\begin{aligned}
\mathcal{L}&=\Bigl[-\frac{p_1}{4\pi}a_\mu\partial_\nu a_\lambda\varepsilon^{\mu\nu\lambda}+\frac{e}{2\pi}A_\mu\partial_\nu a_\lambda\varepsilon^{\mu\nu\lambda}\Bigr]\\
&\quad+\Bigl[-\frac{p_2}{4\pi}\tilde{a}_\mu\partial_\nu\tilde{a}_\lambda\varepsilon^{\mu\nu\lambda}+\frac{1}{2\pi}a_\mu\partial_\nu\tilde{a}_\lambda\varepsilon^{\mu\nu\lambda}\Bigr]
\end{aligned}
\tag{7.3.13}
\end{equation*}
其中 $p_1=m$ 是奇整数。式 (7.3.13) 就是第二级分级 FQH 态的有效理论。

% 块界说明：
% (1) 起点——书页 300（p315）以完整段落收束（"…where p1 = m is an odd integer.
%     Equation (7.3.13) is the effective theory of a second-level hierarchical FQH state."），
%     书页 301 顶部"有效理论可用于确定……"为新段落，故本块无 %--B-TAIL-- 区、正文不以 %JOIN% 起头。
% (2) 终点——书页 313（p328）末句止于 "These charged excitations carry integer charges and"，
%     其续文为 p329（书页 314）顶部 "are created by the electron operators Ψ†.
%     The above theory of edge excitations is formulated in terms of …"。
%     按约定该半句译文（"……携带整数电荷，并且"）留在本文件正文最末，未写入任何 TAIL 区；
%     ch7_d 请以 %JOIN% 起头续译"由电子算符 $\Psi^{\dagger}$ 所产生。……"，勿重复翻译本块已有内容。
% (3) 本块覆盖：7.3.3 后半（自书页 301 顶部起）、7.3.4（含表 7.1、7.2，问题 7.3.3–7.3.6）、
%     7.4 章节首与要点、7.4.1（图 7.15、7.16，脚注 47，问题 7.4.1–7.4.3）、7.4.2 至书页 313 末。
有效理论可以用来确定分级（hierarchical）FQH 态的物理性质。总填充因子由运动方程
$\delta\mathcal{L}/\delta a_{0}\allowbreak=\delta\mathcal{L}/\delta\tilde{a}_{0}\allowbreak=0$ 确定，其中
\begin{equation*}
-eB=p_{1}b-\tilde{b},\qquad b=p_{2}\tilde{b}
\end{equation*}
我们求得
\begin{equation*}
\nu=\frac{b}{-eB}=\frac{1}{p_{1}-\frac{1}{p_{2}}}. \tag{7.3.14}
\end{equation*}
通过引入 $(a_{1\mu},a_{2\mu})=(a_{\mu},\tilde{a}_{\mu})$，式 (7.3.13) 可以写成下面更紧凑的形式：
\begin{equation*}
\mathcal{L}=-\frac{1}{4\pi}K_{IJ}a_{I\mu}\partial_{\nu}a_{J\lambda}\varepsilon^{\mu\nu\lambda}
+\frac{e}{2\pi}q_{I}A_{\mu}\partial_{\nu}a_{I\lambda}\varepsilon^{\mu\nu\lambda} \tag{7.3.15}
\end{equation*}
其中 $K$ 是整数矩阵
\begin{equation*}
K=\begin{pmatrix} p_{1} & -1\\ -1 & p_{2} \end{pmatrix}
\end{equation*}
而 $\boldsymbol{q}$ 是整数矢量 $\boldsymbol{q}^{\top}=(q_{1},q_{2})=(1,0)$。这里 $\boldsymbol{q}$ 将被称为电荷矢量（charge vector）。填充因子 (7.3.14) 可以改写为 $\nu=\boldsymbol{q}^{\top}K^{-1}\boldsymbol{q}$。当 $(p_{1},p_{2})=(3,2)$ 时，该分级态对应于实验中观察到的 $\nu=2/5$ FQH 态。

第二级分级 FQH 态包含两种准粒子。一种是原来电子凝聚体中的准空穴（或涡旋），另一种是新玻色凝聚体中的准空穴（或涡旋）。这两种准空穴分别通过插入源项 $-j^{\mu}a_{\mu}$ 与 $-\tilde{j}^{\mu}\tilde{a}_{\mu}$ 来产生，其中 $\tilde{j}^{\mu}$ 与 $j^{\mu}$ 具有和式 (7.3.2) 类似的形式。第一种准空穴通过在电子波函数上乘以 $\prod_{i}(\xi-z_{i})$ 产生，而第二种准空穴则通过在玻色 Laughlin 波函数上乘以 $\prod_{i}(\eta-\xi_{i})$ 产生（这里 $\xi_{i}$ 是玻色子的复坐标，$\eta$ 是准空穴的位置）。

一个一般的准粒子由 $l_{1}$ 个第一种准粒子和 $l_{2}$ 个第二种准粒子组成，并由两个整数标记。这样的准粒子携带 $l_{1}$ 个单位的 $a_{1\mu}$ 荷和 $l_{2}$ 个单位的 $a_{2\mu}$ 荷，由下式描述
\begin{equation*}
(l_{1}a_{1\mu}+l_{2}a_{2\mu})j^{\mu} \tag{7.3.16}
\end{equation*}
积掉规范场之后，我们发现这样的准粒子携带 $\sum_{J}K_{IJ}l_{J}$ 个单位的 $a_{I\mu}$ 通量。因此，这样的准粒子的统计由下式给出
\begin{equation*}
\theta=\pi\pmb{l}^{\top}K^{-1}\pmb{l}
=\frac{1}{p_{2}p_{1}-1}\left(p_{2}l_{1}^{2}+p_{1}l_{2}^{2}+2l_{1}l_{2}\right)
\end{equation*}
而准粒子的电荷为
\begin{equation*}
Q_{q}=-e\boldsymbol{q}^{\top}K^{-1}\pmb{l}
=-e\frac{p_{2}l_{1}+l_{2}}{p_{2}p_{1}-1}
\end{equation*}
对于 $\nu=2/5$ 态（即 $(p_{1},p_{2})=(3,2)$），具有最小电荷的准粒子由 $(l_{1},l_{2})=(0,1)$ 标记。最小电荷为 $-e/5$。这样的准粒子具有统计 $\theta=\frac{3}{5}\pi$。

我们也可以构造更一般的 FQH 态。这些 FQH 态的有效理论仍具有式 (7.3.15) 给出的形式，只是现在指标 $I$ 从 1 取到某个整数 $n$（$n$ 将被称为该 FQH 态的级（level））。为了得到矩阵 $K$ 的形式，让我们假设在第 $n-1$ 级，有效理论由 $a_{I\mu}$（$I=1,\ldots,n-1$）以及 $K=K^{(n-1)}$ 的式 (7.3.15) 给出。准粒子携带 $a_{I\mu}$ 规范场的整数荷。现在考虑第 $n$ 级分级态，它通过具有 $a_{I\mu}$ 荷 $l_{I}|_{I=1,\ldots,n-1}$ 的准粒子的"凝聚"而得到。这个第 $n$ 级分级态的有效理论将由含 $n$ 个规范场的式 (7.3.15) 给出。第 $n$ 个规范场 $a_{n\mu}$ 来自新的凝聚体。矩阵 $K$ 由下式给出
\begin{equation*}
K^{(n)}=\begin{pmatrix} K^{(n-1)} & -\pmb{l}\\ -\pmb{l}^{\top} & p_{n} \end{pmatrix}
\end{equation*}
其中 $p_{n}$ 为偶数。电荷矢量 $\boldsymbol{q}$ 仍由 $(1,0,0,\ldots)$ 给出。通过迭代我们看到，广义分级态总是由整数对称矩阵描述，其中 $K_{II}$ 为偶数，唯独 $K_{11}$ 为奇数。新的凝聚体给出一种新的准粒子，它同样携带新规范场 $a_{n\mu}$ 的整数荷。因此，一个一般的准粒子总是携带 $a_{I\mu}$ 场的整数荷。

让我们用更一般的语言总结上述结果。我们知道，一个分级（或广义分级）FQH 态包含许多不同的凝聚体。这些凝聚体并不相互独立：一个凝聚体中的粒子对其他凝聚体中的粒子而言就像一根通量管（flux tube）。为了描述这种耦合，方便的做法是用 $U(1)$ 规范场来描述第 $I$ 个凝聚体的密度和流 $J_{I\mu}$：
\begin{equation*}
J_{I}^{\mu}=\frac{1}{2\pi}\varepsilon^{\mu\alpha\beta}\partial_{\alpha}a_{I\beta} \tag{7.3.17}
\end{equation*}
这时，不同凝聚体之间的耦合由规范场的 Chern--Simons 项描述：
\begin{equation*}
\mathcal{L}=-\frac{1}{4\pi}K_{IJ}a_{I\mu}\partial_{\nu}a_{J\lambda}\varepsilon^{\mu\nu\lambda}
+\frac{e}{2\pi}q_{I}A_{\mu}\partial_{\nu}a_{I\lambda}\varepsilon^{\mu\nu\lambda} \tag{7.3.18}
\end{equation*}
当 $K$ 取为一般的整数 $n\times n$ 矩阵，且 $K_{II}|_{I=1}$ 为奇数、$K_{II}|_{I>1}$ 为偶数、$\boldsymbol{q}^{\top}=(1,0,\ldots,0)$ 时，式 (7.3.18) 描述的是最一般的（阿贝尔）FQH 态 (Wen and Zee, 1992a)。填充因子由下式给出
\begin{equation*}
\nu=\boldsymbol{q}^{\top}K^{-1}\boldsymbol{q} \tag{7.3.19}
\end{equation*}
准粒子激发可以看作不同凝聚体中的涡旋。一个一般的准粒子由 $n$ 个整数 $l_{I}|_{I=1,..,n}$ 标记，并可由下面的源项产生：
\begin{equation*}
\mathcal{L}=l_{I}a_{I\mu}j^{\mu} \tag{7.3.20}
\end{equation*}
这样的准粒子将记作 $\psi_{\pmb{l}}$。式 (7.3.20) 中的 $j^{\mu}$ 具有如下形式
\begin{gather*}
\pmb{j}(x)=\dot{\pmb{x}}_{0}\delta(\pmb{x}-\pmb{x}_{0})\\
j^{0}(x)=\delta(\pmb{x}-\pmb{x}_{0}) \tag{7.3.21}
\end{gather*}
它在 $\pmb{x}_{0}$ 处产生一个准粒子。

当我们产生一个准粒子 $\psi_{\pmb{l}}$ 时，它会在所有凝聚体的密度中引起一个改变 $\delta J_{I}^{0}$。由式 (7.3.18) 与式 (7.3.20) 的运动方程，我们发现 $\delta J_{I}^{0}$ 满足
\begin{equation*}
\int\mathrm{d}^{2}x\,\delta J_{I}^{0}
=l_{J}(K^{-1})_{JI}\int\mathrm{d}^{2}x\,j^{0}=(\pmb{l}^{\top}K^{-1})_{I} \tag{7.3.22}
\end{equation*}
准粒子 $\psi_{\pmb{l}}$ 的电荷与统计由下式给出
\begin{equation*}
\theta_{\pmb{l}}=\pi\pmb{l}^{\top}K^{-1}\pmb{l},\qquad
Q_{\pmb{l}}=-e q_{I}\int\mathrm{d}^{2}x\,\delta J_{I}^{0}=-e\pmb{l}^{\top}K^{-1}\boldsymbol{q} \tag{7.3.23}
\end{equation*}
准粒子统计的这一结果来自如下观察：$2\pi\delta J_{I}^{0}$ 是 $a_{I\mu}$ 的通量，而准粒子携带 $l_{I}$ 个单位的 $a_{I\mu}$ 荷。

一个一般的电子激发可以看作一种特殊的（一般）准粒子，$\psi_{e}=\psi_{\pmb{l}_{e}}$，其中整数矢量 $\pmb{l}_{e}$ 由下式给出
\begin{equation*}
l_{eI}=K_{IJ}L_{J},\qquad L_{I}=\text{整数},\qquad q_{I}L_{I}=1 \tag{7.3.24}
\end{equation*}
我们可以证明，这些电子激发满足下列性质：它们携带单位电荷（见式 (7.3.23)）；它们具有费米统计；让一个电子激发 $\psi_{e}=\psi_{L}$ 绕任意准粒子激发 $\psi_{\pmb{l}}$ 一周，总是诱导 $2\pi$ 整数倍的相位；并且由式 (7.3.24) 定义的激发，就是满足上述三条性质的全部激发。

由式 (7.3.15) 的一般有效理论，我们发现广义分级态可以由一个整数值的 $K$ 矩阵和一个电荷矢量 $\boldsymbol{q}$ 来标记。现在我们想问这样一个问题：不同的 $(K,\boldsymbol{q})$ 是否描述不同的 FQH 态？注意到，通过对规范场 $a_{I\mu}$ 的重新定义，总可以把 $K$ 矩阵对角化成对角元为 $\pm1$ 的矩阵。这样看来，所有具有相同符号差（signature）的 $K$ 矩阵似乎都描述相同的 FQH 态，因为经过规范场的适当重新定义后，它们给出同一个有效理论。这个结论当然是错误的。我们想强调，仅凭有效拉格朗日量 (7.3.15) 并不能恰当描述分级态的内部序（或拓扑序）。正是有效拉格朗日量 (7.3.15) 连同 $a_{I\mu}$ 荷的量子化条件一起，才刻画了拓扑序。配备了允许 $U(1)$ 荷之量子化条件的 $U(1)$ 规范理论被称为紧致 $U(1)$ 理论。我们的有效理论 (7.3.15) 实际上是一个所有 $U(1)$ 荷都量子化为整数的紧致 $U(1)$ 理论。因此，允许的 $U(1)$ 荷构成一个 $n$ 维立方格子，我们将称之为电荷格子（charge lattice）。所以，当考虑两个不同 $K$ 矩阵的等价性时，只能使用保持电荷量子化条件不变（即保持电荷格子不变）的场重新定义。把电荷格子映到其自身的变换属于群 $SL(n,Z)$（行列式为 1 的整数矩阵群）：
\begin{equation*}
a_{I\mu}\rightarrow W_{IJ}a_{J\mu},\qquad W\in SL(n,Z) \tag{7.3.25}
\end{equation*}
由以上讨论我们看到，由 $(K_{1},\boldsymbol{q}_{1})$ 与 $(K_{2},\boldsymbol{q}_{2})$ 描述的两个 FQH 态是等价的（即它们属于同一个普适类（universality class）），如果存在 $W\in SL(n,Z)$ 使得
\begin{equation*}
\boldsymbol{q}_{2}=W\boldsymbol{q}_{1},\qquad K_{2}=WK_{1}W^{\top} \tag{7.3.26}
\end{equation*}
这是因为在变换 (7.3.25) 之下，由 $(K_{1},\boldsymbol{q}_{1})$ 描述的有效理论简单地变成由 $(K_{2},\boldsymbol{q}_{2})$ 描述的另一个有效理论。

我们想指出，在上述讨论中我们忽略了另一个拓扑量子数——自旋矢量（spin vector）。因此，式 (7.3.26) 的等价条件不适用于旋转不变系统。但是，式 (7.3.26) 确实适用于无序 FQH 系统，因为在无序系统中角动量不守恒，自旋矢量没有良好定义。关于自旋矢量的讨论可参见 Wen and Zee (1992c,d)。表 7.1 列出了一些常见的单层自旋极化 FQH 态的 $K$ 矩阵、电荷矢量 $\boldsymbol{q}$ 与自旋矢量 $\boldsymbol{s}$。

\subsection*{问题 7.3.3}
证明式 (7.3.24) 之下列出的四条性质。

\subsection{简单多层的分数量子霍尔态的有效理论}

\begin{keypoints}
\item $K$ 矩阵与多层 FQH 波函数。
\end{keypoints}

\begin{center}
\small 表 7.1\quad 某些简单的单层自旋极化 FQH 态的 $(K,\boldsymbol{q},\boldsymbol{s})$ 表达式\\[6pt]
\begin{tabular}{cccc}
\toprule
$\nu$ & $\boldsymbol{q}$ & $K$ & $\boldsymbol{s}$ \\
\midrule
$1/m$ & $(1)$ & $(m)$ & $\begin{pmatrix} m/2 \end{pmatrix}$ \\[10pt]
$1-1/m$ & $\begin{pmatrix}1\\0\end{pmatrix}$ &
$\begin{pmatrix}1 & 1\\ 1 & -(m-1)\end{pmatrix}$ &
$\begin{pmatrix}1/2\\ (1-m)/2\end{pmatrix}$ \\[18pt]
$2/5$ & $\begin{pmatrix}1\\0\end{pmatrix}$ &
$\begin{pmatrix}3 & -1\\ -1 & 2\end{pmatrix}$ &
$\begin{pmatrix}1/2\\ 1\end{pmatrix}$ \\[18pt]
$3/7$ & $\begin{pmatrix}1\\0\\0\end{pmatrix}$ &
$\begin{pmatrix}3 & -1 & 0\\ -1 & 2 & -1\\ 0 & -1 & 2\end{pmatrix}$ &
$\begin{pmatrix}1/2\\ 1\\ 1\end{pmatrix}$ \\
\bottomrule
\end{tabular}
\end{center}

到目前为止，我们只考虑了所谓的单层 FQH 态，其中二维电子气位于两种不同半导体的界面上。我们也可以制作具有多层界面的样品。置于强磁场之下时，不同层上的二维电子气可以形成多层 FQH 态。

用于构造分级态有效理论的同一方法，也可用来构造多层 FQH 态的有效理论。对多层 FQH 态而言，FQH 波函数与 $K$ 矩阵之间的联系变得非常透明。本节将集中讨论双层 FQH 态，不过推广到 $n$ 层 FQH 态是直截了当的。

我们想为下面这个简单的双层 FQH 态构造有效理论：
\begin{equation*}
\prod_{i<j}(z_{1i}-z_{1j})^{l}\prod_{i<j}(z_{2i}-z_{2j})^{m}\prod_{i,j}(z_{1i}-z_{2j})^{n}\,
\mathrm{e}^{-\frac{1}{4l_{B}^{2}}\left(\sum_{i}|z_{1i}|^{2}+\sum_{j}|z_{2j}|^{2}\right)}, \tag{7.3.27}
\end{equation*}
其中 $z_{Ii}$ 是第 $I$ 层中第 $i$ 个电子的复坐标。这里 $l$ 和 $m$ 是奇整数，使波函数与电子的费米统计相一致；$n$ 则可以是任意非负整数。上述波函数由 Halperin 首先提出，作为 Laughlin 波函数的推广 (Halperin, 1983)。这些波函数似乎能够解释双层 FQH 系统中观测到的部分主要 FQH 填充因子。

我们从第一层的单层 FQH 态出发，即 $\prod_{i<j}(z_{1i}-z_{1j})^{l}\mathrm{e}^{-\sum_{i}|z_{1i}|^{2}/4}$，它是一个 $1/l$ Laughlin 态，由如下有效理论描述
\begin{equation*}
\mathcal{L}=\left[-l\frac{1}{4\pi}a_{1\mu}\partial_{\nu}a_{1\lambda}\varepsilon^{\mu\nu\lambda}
+\frac{e}{2\pi}A_{\mu}\partial_{\nu}a_{\lambda}\varepsilon^{\mu\nu\lambda}\right] \tag{7.3.28}
\end{equation*}
其中 $a_{1\mu}$ 是描述第一层中电子密度和流的规范场。

考察式 (7.3.27) 的波函数可以看到，第二层中的一个电子束缚于第一层中的一个准空穴激发。这样的准空穴激发由 $n$ 个基本准空穴激发组成，并携带 $a_{1\mu}$ 荷 $-n$。一团准空穴由下面的有效理论描述：
\begin{equation*}
\mathcal{L}=-na_{1\mu}j^{\mu}+\text{动能/势能} \tag{7.3.29}
\end{equation*}
其中 $j^{\mu}$ 具有式 (7.3.2) 给出的形式。如前所述，在平均场理论中，如果忽略 $a_{1\mu}$ 场的响应，则式 (7.3.29) 简单地描述"磁场" $nb_{1}$ 中的一团玻色子，这里 $b_{1}=-\varepsilon_{ij}\partial_{i}a_{1j}$。现在我们想给式 (7.3.29) 中的每个准空穴附着一个（第二层中的）电子。这一操作有以下两个效果：由于电子携带电荷 $-e$，准空穴与电子的束缚态可以直接与电磁场 $A_{\mu}$ 耦合；并且束缚态表现得像一个费米子。束缚态的有效理论具有如下形式：
\begin{equation*}
\mathcal{L}=(eA_{\mu}-na_{1\mu})j^{\mu}+\text{动能/势能} \tag{7.3.30}
\end{equation*}
它现在（在平均场理论中）描述一团费米子。这些费米子感受到一个有效磁场 $-eB+nb_{1}$。

当第二层中的电子（即式 (7.3.30) 中的费米子）具有密度 $\frac{1}{m}(-eB+nb_{1})/2\pi$ 时（即它们具有有效填充因子 $1/m$），它们可以形成一个 $1/m$ Laughlin 态，对应于波函数中的 $\prod_{i<j}(z_{2i}-z_{2j})^{m}$ 部分。（注意这里 $eB<0$。）引入新规范场 $j^{\mu}=\frac{1}{2\pi}\partial_{\nu}a_{2\lambda}\varepsilon^{\mu\nu\lambda}$ 来描述式 (7.3.30) 中的费米子流 $j^{\mu}$，第二层中 $1/m$ 态的有效理论便具有形式
\begin{equation*}
\mathcal{L}=-m\frac{1}{4\pi}a_{2\mu}\partial_{\nu}a_{2\lambda}\varepsilon^{\mu\nu\lambda} \tag{7.3.31}
\end{equation*}
把式 (7.3.28)、式 (7.3.30) 与式 (7.3.31) 合在一起，我们就得到双层态的总有效理论，它具有式 (7.3.15) 给出的形式，其中 $K$ 矩阵与电荷矢量 $\boldsymbol{q}$ 由下式给出
\begin{equation*}
K=\begin{pmatrix} l & n\\ n & m \end{pmatrix},\qquad
\boldsymbol{q}=\begin{pmatrix}1\\1\end{pmatrix} \tag{7.3.32}
\end{equation*}
我们看到，$K$ 矩阵的矩阵元就是波函数中的幂次。该 FQH 态的填充因子仍由式 (7.3.19) 给出。

双层态中有两种（基本）准空穴激发。第一种通过在基态波函数上乘以 $\prod_{i}(\xi-z_{1i})$ 产生，第二种则由 $\prod(\xi-z_{2i})$ 产生。作为第一层与第二层这两个电子凝聚体中的涡旋，第一种准空穴由源项 $-a_{1\mu}j^{\mu}$ 产生，第二种准空穴由 $-a_{2\mu}j^{\mu}$ 产生。因此，双层态中一个一般的准粒子是若干第一类与第二类准空穴的束缚态，由式 (7.3.16) 描述。这些准粒子的量子数仍由式 (7.3.23) 给出。

\begin{center}
\small 表 7.2\quad 某些双层 FQH 态的 $(K,\boldsymbol{q},\boldsymbol{s})$ 表达式\\[6pt]
\begin{tabular}{cccc}
\toprule
$\nu$ & $\boldsymbol{q}$ & $K$ & $\boldsymbol{s}$ \\
\midrule
$1/m$ & $\begin{pmatrix}1\\1\end{pmatrix}$ &
$\begin{pmatrix}m & m\\ m & m\end{pmatrix}$ &
$\begin{pmatrix}m/2\\ m/2\end{pmatrix}$ \\[18pt]
$2/5$ & $\begin{pmatrix}1\\1\end{pmatrix}$ &
$\begin{pmatrix}3 & 2\\ 2 & 3\end{pmatrix}$ &
$\begin{pmatrix}3/2\\ 3/2\end{pmatrix}$ \\[18pt]
$1/2$ & $\begin{pmatrix}1\\1\end{pmatrix}$ &
$\begin{pmatrix}3 & 1\\ 1 & 3\end{pmatrix}$ &
$\begin{pmatrix}3/2\\ 3/2\end{pmatrix}$ \\[18pt]
$2/3$ & $\begin{pmatrix}1\\1\end{pmatrix}$ &
$\begin{pmatrix}3 & 0\\ 0 & 3\end{pmatrix}$ &
$\begin{pmatrix}3/2\\ 3/2\end{pmatrix}$ \\[18pt]
$2/3$ & $\begin{pmatrix}1\\1\end{pmatrix}$ &
$\begin{pmatrix}1 & 2\\ 2 & 1\end{pmatrix}$ &
$\begin{pmatrix}1/2\\ 1/2\end{pmatrix}$ \\
\bottomrule
\end{tabular}
\end{center}

一般地，(7.3.27) 类型的多层 FQH 态由一个矩阵元为整数、对角元为奇整数的 $K$ 矩阵描述。电荷矢量具有形式 $\boldsymbol{q}^{\top}=(1,1,\ldots,1)$。人们通常把双层 FQH 态 (7.3.27) 标记为 $(l,m,n)$。表 7.2 列出了一些简单双层态的 $K$ 矩阵、电荷矢量 $\boldsymbol{q}$ 与自旋矢量 $\boldsymbol{s}$。

由上述我们看到，(332) 双层态具有填充因子 2/5——这个填充因子也出现在单层分级态中。于是产生如下问题：双层 2/5 态与单层 2/5 态是否属于同一普适类？这个问题具有实验上的后果。我们可以设想这样的实验：从一个层间隧穿非常弱的系统中的 (332) 双层态出发，在保持填充因子不变的同时，把层间隧穿变得越来越强，双层态最终将变成单层 2/5 态。问题是：双层 (332) 态与单层 2/5 态之间的转变是平滑的渡越（crossover）还是相变？如果忽略自旋矢量，我们看到两个 2/5 态的 $K$ 矩阵和电荷矢量是等价的，因为它们通过一个 $SL(2,Z)$ 变换相联系。因此，在没有旋转对称性的情形下（此时自旋矢量没有良好定义），两个 2/5 态可以平滑地彼此转变。而当我们把自旋矢量包括进来时，两个 2/5 态并不等价，并且对旋转不变系统而言，它们被一个一级相变分隔开。

从表 7.2 我们还看到，存在两种不同的双层 2/3 态。当层内相互作用远强于层间相互作用时（真实样品中正是这种情形），基态波函数倾向于让同层电子之间具有更高阶的零点。因此，(330) 态的能量应当比 (112) 态低。单层 2/3 态的 $K$ 矩阵和电荷矢量与 (112) 态的等价，而与 (330) 态的不等价。因此，要把一个双层 2/3 态（即 (330) 态）变成单层 2/3 态，无论旋转对称性如何都必须经过一次相变。

双层 (mmm) 态是一个有趣的态，因为 $\det K=0$。Fertig (1989)、Brey (1990) 与 MacDonald 等 (1990) 研究了 (111) 态，发现了一个无能隙的集体模式。Wen 和 Zee (1992b) 研究了更一般的 (mmm) 态，并指出 (mmm) 态（在无层间隧穿时）自发破缺一个 $U(1)$ 对称性，是一种中性超流体。更详细的讨论（包括其实验含义）可参见 Wen and Zee (1993)、Murphy et al. (1994) 与 Yang et al. (1994)。

\subsection*{问题 7.3.4}
让我们首先忽略自旋矢量 $\boldsymbol{s}$。证明表 7.1 中单层 2/5 态的 $(K,\boldsymbol{q})$ 与表 7.2 中双层 2/5 态的 $(K,\boldsymbol{q})$ 等价。再证明，这两个态的 $(K,\boldsymbol{q},\boldsymbol{s})$ 却不等价。（在存在自旋矢量 $\boldsymbol{s}$ 时，等价关系变为 $\boldsymbol{s}_{2}=W\boldsymbol{s}_{1}$、$\boldsymbol{q}_{2}=W\boldsymbol{q}_{1}$ 与 $K_{2}=WK_{1}W^{+}$。）

\subsection*{问题 7.3.5}
考虑双层 (lmn) 态 (7.2.6)。利用有效理论 (7.3.15) 与式 (7.3.32) 求出填充因子。并求出第一层与第二层中准空穴的分数电荷与分数统计。（可以把结果与问题 7.2.5 的结果比较。）

\subsection*{问题 7.3.6}
假设双层 (mmm) 态的有效理论具有形式
\begin{equation*}
\mathcal{L}=-\frac{1}{4\pi}K_{IJ}a_{I\mu}\partial_{\nu}a_{J\lambda}\varepsilon^{\mu\nu\lambda}
+\frac{1}{2g_{1}}\pmb{e}_{I}\cdot\pmb{e}_{I}-\frac{1}{2g_{2}}b_{I}b_{I}
\end{equation*}
其中 $\pmb{e}_{I}$ 与 $b_{I}$ 分别是 $a_{I\mu}$ 的"电场"与"磁场"。
\begin{enumerate}
\item 证明 (mmm) 态正如超流体一样具有无能隙激发。求出无能隙激发的速度。
\item 同样像超流体一样，(mmm) 态包含能量为 $\gamma\ln L$ 的涡旋激发，这里 $L$ 是系统的线度。求出能量最小的涡旋的 $\gamma$ 值。（提示：在 (mmm) 态中，$l_{I}$ 的 $a_{I\mu}$ 荷仍量子化为整数。）
\item 如果我们确实把 (mmm) 态当作超流体，你能辨认出破缺的对称性和序参量吗？
\end{enumerate}

\section{分数量子霍尔液体中的边缘激发}

\begin{keypoints}
\item FQH 态\emph{总是}具有无能隙的边缘激发。
\item FQH 态的边缘激发构成手征 Luttinger 液体。
\item 边缘激发的结构由体拓扑序决定。
\end{keypoints}

由于电子之间排斥性的相互作用和强关联，尽管第一朗道能级只是部分填充，QH 液体仍是一个不可压缩态。QH 态中的所有体激发都具有有限的能隙。QH 态与绝缘体非常相似：两者都具有有限能隙和短程的电子传播子。正是由于这种相似性，人们曾对如下事实感到困惑：QH 系统的输运性质显然与普通绝缘体截然不同。Halperin 首先指出，IQH 态含有无能隙的边缘激发 (Halperin, 1982)。IQH 态的非平庸输运性质正来自这些无能隙边缘激发 (Halperin, 1982; Trugman, 1983; MacDonald and Streda, 1984; Streda et al., 1987; Buttiker, 1988; Jain and Kivelson, 1988a,b)。例如，对一个 IQH 样品做两端测量，只有当源与汇通过边缘连通时才可能给出有限电阻。如果源与汇不被任何边缘连通，那么两端测量在零温下将给出无穷大电阻——这是一个与绝缘体十分相似的结果。Halperin 还研究了 IQH 态边缘激发的动力学性质，发现边缘激发由一个手征的一维费米液体理论描述。

鉴于 FQH 态与 IQH 态输运性质的相似性，一个自然的猜想是：FQH 态的输运同样由边缘激发支配 (Beenakker, 1990; MacDonald, 1990)。然而，由于 FQH 态本质上是多体态，FQH 态中的边缘激发无法通过填充单粒子能级来构造。换言之，FQH 态的边缘激发不应当用费米液体来描述。因此，我们需要一种全新的方法来理解 FQH 液体边缘态的动力学性质 (Wen, 1992, 1995)。

在下一节中，我们将用费米液体理论研究 $\nu=1$ IQH 态的边缘激发。然后，我们将利用流代数（或玻色化，见 Tomonaga, 1950; Luttinger, 1963; Coleman, 1975）来建立 FQH 态边缘激发的理论。

\subsection{整数量子霍尔边缘态的费米液体理论}

\begin{keypoints}
\item IQH 液体的基态通过填充单粒子能级得到。因此，低能边缘激发可以通过以略有不同的方式填充这些单粒子能级来获得。
\end{keypoints}

考虑均匀磁场 $B$ 中的一团无相互作用电子气。在对称规范下，角动量是一个好量子数。能量与动量的共同本征态构成一个环状分布（见图 7.10）。因此，在光滑的圆形势 $V(r)$ 存在时，单粒子能级如图 7.15 所示。基态通过填充能量最低的 $N$ 个能级得到。这样的态对应于一个具有均匀密度的圆形液滴（见图 7.11）。

\begin{figure}[H]
\centering
\includegraphics[width=0.72\textwidth]{fig_7.15.pdf}
\caption*{图 7.15\quad 光滑势 $V(r)$ 存在时头三个朗道能级中的能级。这里 $m$ 是能级的角动量。化学势 $\mu$ 以下的能级各被一个电子占据，从而给出 IQH 态。改变边缘附近的占据数，便产生 IQH 态的无能隙边缘激发。}
\end{figure}

在零朗道能级中，角动量 $m$ 的态的能量为
\begin{equation*}
E_{m}=\frac{1}{2}\hbar\omega_{c}+V(r_{m})
\end{equation*}
其中 $r_{m}=\sqrt{2m}\,l_{B}$ 是 $m$ 态的半径。从图 7.15 我们看到，体激发具有有限的能隙 $\hbar\omega_{c}$。边缘激发可以这样产生：把一个电子从边缘附近的 $m$ 态移到 $m+1$ 态（见图 7.15）。由于在热力学极限 $r_{m}\to\infty$ 下 $r_{m+1}-r_{m}\to0$，边缘激发是无能隙的。

由于 $m$ 态具有半径 $r_{m}$，我们也可以把角动量 $m$ 看作沿边缘的动量 $k=m/r_{m}$。这使我们能够把 $E_{m}$ 当作一个能量--动量关系
\begin{equation*}
E(k)=\frac{1}{2}\hbar\omega_{c}+V(2l_{B}^{2}k)
\end{equation*}
用 $E(k)$ 来表示，我们的电子系统可以用一个无相互作用的哈密顿量描述：
\begin{equation*}
H=\sum_{k}E(k)c_{k}^{\dagger}c_{k} \tag{7.4.1}
\end{equation*}
上述哈密顿量描述一个一维手征费米液体（chiral Fermi liquid）。我们称它为手征费米液体，是因为它只有一个费米点，并且所有低能激发都朝同一个方向传播。\footnote{相比之下，通常的一维费米液体有两个费米点，同时包含左移与右移两种成分。}我们由此得出结论：$\nu=1$ IQH 态的边缘激发由一维手征费米液体 (7.4.1) 描述。

\begin{figure}[H]
\centering
\includegraphics[width=0.72\textwidth]{fig_7.16.pdf}
\caption*{图 7.16\quad (a) 能量--角动量关系 $E_{m}$ 可以看作能量--动量关系 $E(k)$。边缘激发可以看作一维手征费米液体的激发，其中只含右移成分。(b) 通常的一维费米液体的色散关系同时包含右移与左移成分。}
\end{figure}

\subsection*{问题 7.4.1}
证明 IQH 边缘激发的速度 $v=\partial E(k)/\partial k$ 由 $cE/B$ 给出，其中 $c$ 是光速，$E=-\partial V(r)/\partial r$ 是约束势 $V(r)$ 在边缘处产生的电场。设 $\nu=1$ IQH 液滴包含 $N$ 个电子，求出手征费米液体的费米动量 $k_{F}$。

\subsection*{问题 7.4.2}
一个 $\nu=1$ IQH 态被圆形光滑势 $V(r)$ 约束，包含 $N$ 个电子。求出基态的总角动量 $M_{0}$。求出总角动量为 $M_{0}+m$ 的低能粒子--空穴激发的数目，其中 $m=-1$ 及 $m=1,2,3,4,5$。证明在热力学极限下，具有相同总角动量的激发具有相同的能量。

\subsection*{问题 7.4.3}
考虑一个由 $H=\sum_{k}(v_{1}k\,c_{1,k}^{\dagger}c_{1,k}+v_{2}k\,c_{2,k}^{\dagger}c_{2,k})$ 描述的一维无相互作用费米子系统。混合项 $\sum_{k}(\gamma c_{1,k}^{\dagger}c_{2,k}+\text{h.c.})$ 是一个相关微扰。证明：当系统同时具有右移与左移成分（即 $v_{1}v_{2}<0$）时，这个相关微扰会打开一个有限的能隙，并剧烈地改变系统的低能性质。然而，对于所有激发都朝同一方向运动的手征费米液体（即 $v_{1}v_{2}>0$），这个相关微扰不诱导任何能隙。事实上，我们认为任何微扰都无法在一维手征费米液体中诱导出能隙。一维手征费米液体中的无能隙激发对所有微扰都是稳健的，并且在拓扑上是稳定的。

\subsection{流体动力学方法——$1/m$ Laughlin 态}

\begin{keypoints}
\item 边缘激发来自流代数（Kac--Moody 代数）。
\item 无法通过填充单粒子能级来构造 FQH 边缘激发。
\end{keypoints}

由于 $1/m$ Laughlin FQH 态不能通过填充单粒子能级得到，我们不能用上一节的图像来构造 $1/m$ Laughlin 态的边缘激发。由于不知道如何从相互作用电子的哈密顿量导出 FQH 边缘理论，下面我们转而尝试猜测一个低能有效理论。

猜测边缘激发动力学理论的最简单方法是使用流体动力学方法。在这种方法中，人们利用这样一个事实：FQH 态是不可压缩、无旋的液体，不含低能体激发。因此，唯一的低能激发（在体能隙之下）是 FQH 液滴形变产生的表面波。这些表面波被认定为 FQH 态的边缘激发。

在流体动力学方法 (Wen, 1992) 中，我们首先研究 FQH 液滴上表面波的经典理论，然后把经典理论量子化，得到边缘激发的量子描述。令人惊叹的是，由经典理论得到的这个简单的量子描述，在低能下给出了边缘激发的完整描述，并使我们能够计算沿边缘的电子与准粒子传播子。

考虑一个填充因子为 $\nu$、被势阱约束的 FQH 液滴。由于电导非零，势阱的电场产生一个沿边缘流动的持续电流：
\begin{equation*}
\pmb{j}=\sigma_{xy}\hat{z}\times\pmb{E},\qquad \sigma_{xy}=\nu\frac{e^{2}}{h}
\end{equation*}
这意味着边缘附近的电子以速度
\begin{equation*}
v=\frac{E}{B}c
\end{equation*}
沿边缘漂移，这里 $c$ 是光速。于是，边缘波（边缘的形变）也以速度 $v$ 传播。让我们用一维密度 $\rho(x)=nh(x)$ 来描述边缘波，其中 $h(x)$ 是边缘的位移，$x$ 是沿边缘的坐标，而 $n=\nu/2\pi l_{B}^{2}$ 是体内的二维电子密度。（这里 $l_{B}=\sqrt{c/eB}$ 是磁长度。）我们看到，边缘波的传播由下面的波动方程描述：
\begin{equation*}
\partial_{t}\rho+v\partial_{x}\rho=0 \tag{7.4.2}
\end{equation*}
注意，边缘波总是朝一个方向传播；不存在朝相反方向传播的波。

边缘波的哈密顿量（即能量）由下式给出
\begin{equation*}
H=\int\mathrm{d}x\,\frac{1}{2}e\rho E h=\int\mathrm{d}x\,\pi\frac{v}{\nu}\rho^{2} \tag{7.4.3}
\end{equation*}
其中我们用到了 $h=\rho/n=\rho\frac{2\pi c}{eB}$。在动量空间中，式 (7.4.2) 与式 (7.4.3) 可以改写为
\begin{equation*}
\dot{\rho}_{k}=-\mathrm{i}vk\rho_{k},\qquad
H=2\pi\frac{v}{\nu}\sum_{k>0}\rho_{-k}\rho_{k} \tag{7.4.4}
\end{equation*}
其中 $\rho_{k}=\int\mathrm{d}x\,\frac{1}{\sqrt{L}}\mathrm{e}^{\mathrm{i}kx}\rho(x)$，而 $L$ 是边缘的长度。我们发现，如果把 $\rho_{k}|_{k>0}$ 认作"坐标"，并把 $\pi_{k}=\mathrm{i}2\pi\rho_{-k}/\nu k$ 认作相应的正则"动量"，那么标准哈密顿方程
\begin{equation*}
\dot{q}=\frac{\partial H}{\partial p},\qquad \dot{p}=-\frac{\partial H}{\partial q}
\end{equation*}
将重现运动方程 $\dot{\rho}_{k}=\mathrm{i}vk\rho_{k}$。这使我们得以认定正则"坐标"与"动量"。有趣的是，位移 $h(x)$ 同时包含"坐标"与"动量"。这是边缘波的手征性质所致。

知道了正则坐标与动量之后，量子化经典理论就很容易了。我们只需把 $\rho_{k}$ 和 $\pi_{k}$ 看作满足 $[\rho_{k},\pi_{k'}]=\mathrm{i}\delta_{kk'}$ 的算符。于是，量子化之后我们得到
\begin{align*}
[\rho_{k},\rho_{k'}]&=\frac{\nu}{2\pi}k\delta_{k+k'},\qquad
k,k'=\text{整数}\times\frac{2\pi}{L} \tag{7.4.5}\\
H&=2\pi\frac{v}{\nu}\sum_{k>0}\rho_{-k}\rho_{k}
\end{align*}
上述代数称为 (U(1)) Kac--Moody（K--M）代数 (Kac, 1983; Goddard and Olive, 1985, 1986)。类似的代数也曾出现在 Tomonaga 模型中 (Tomonaga, 1950; Luttinger, 1963; Coleman, 1975)。注意，式 (7.4.5) 简单地描述了一组退耦谐振子（由 $(\rho_{k},\rho_{-k})$ 生成）。因此，式 (7.4.5) 是一个一维自由声子理论（只有单一分支）。对于 $k>0$，$\rho_{k}^{\dagger}=\rho_{-k}$ 产生一个动量为 $k$、能量为 $vk$ 的声子，而 $\rho_{k}$ 湮灭一个声子。式 (7.4.5) 为 Laughlin 态的低能边缘激发提供了完整的描述。

这里考虑的边缘激发不改变系统的总电荷，因而是中性的。下面我们将讨论带电激发，并由 K--M 代数 (7.4.5) 计算电子传播子。

低能电荷激发显然对应于把电子加到边缘上（从边缘移走）。这些带电激发携带整数电荷，并且
由电子算符 $\Psi^\dagger$ 产生。上述边缘激发理论是用一维密度算符 $\rho(x)$ 来表述的。因此，中心任务是把电子算符用密度算符表示出来。边缘上的电子算符产生一个局域电荷，故应满足
\begin{equation*}
[\rho(x),\Psi^\dagger(x')]=\delta(x-x')\Psi^\dagger(x')
\tag{7.4.6}
\end{equation*}
由于 $\rho$ 满足 K--M 代数 (7.4.5)，可以证明
\begin{equation*}
[\rho(x_1),\rho(x_2)]=\mathrm{i}\frac{\nu}{2\pi}\delta'(x_1-x_2)
\end{equation*}
或者
\begin{equation*}
[\rho(x_1),\phi(x_2)]=-\mathrm{i}\nu\delta(x_1-x_2)
\end{equation*}
其中 $\phi$ 由 $\rho=\frac{1}{2\pi}\partial_x\phi$ 给出。我们看到，$\rho(x)$ 可以视为 $\phi$ 的泛函导数，即 $\rho(x)=-\mathrm{i}\nu\frac{\delta}{\delta\phi(x)}$。利用这一关系，可以证明满足式 (7.4.6) 的算符由下式给出（Wen, 1992）：
\begin{equation*}
\Psi\propto\mathrm{e}^{\mathrm{i}\frac{1}{\nu}\phi}
\tag{7.4.7}
\end{equation*}

式 (7.4.6) 只表明算符 $\Psi$ 携带电荷 $e$。为了把 $\Psi$ 认作电子算符，我们还需要证明 $\Psi$ 是费米性算符。利用 K--M 代数 (7.4.5)，我们发现（见问题 7.4.8）
\begin{equation*}
\Psi(x)\Psi(x')=(-)^{1/\nu}\Psi(x')\Psi(x)
\tag{7.4.8}
\end{equation*}

我们看到，式 (7.4.7) 中的电子算符 $\Psi$ 只有当 $1/\nu=m$ 为奇整数时才是费米性的，此时该 FQH 态是一个 Laughlin 态。

在上述讨论中，我们做了一个并不普遍成立的假设。我们假设了不可压缩 FQH 液体只包含\emph{一种}不可压缩流体组分，从而只导致一支边缘激发。上述结果意味着，当 $\nu\neq1/m$ 时，只有一支的边缘理论不包含电子算符，因而是不自洽的。我们的单支边缘理论只适用于 $\nu=1/m$ 的 Laughlin 态。

现在让我们计算 $\nu=1/m$ Laughlin 态边缘上的电子传播子。由于 $\phi$ 是一个自由声子场，其传播子为
\begin{equation*}
\langle\phi(x,t)\phi(0)\rangle=-\nu\ln(x-vt)+\text{常数}
\end{equation*}
电子传播子可以容易地算出为
\begin{equation*}
G(x,t)=\langle T(\Psi^\dagger(x,t)\Psi(0))\rangle=\exp[\frac{1}{\nu^2}\langle\phi(x,t)\phi(0)\rangle]\propto\frac{1}{(x-vt)^m}
\tag{7.4.9}
\end{equation*}

我们看到的第一件事是：FQH 态边缘上的电子传播子获得了一个不等于 1 的非平庸指数 $m=1/\nu$。这意味着 FQH 态边缘上的电子是强关联的，不能用费米液体理论来描述。我们将把这类电子态称为手征 Luttinger 液体（chiral Luttinger liquid）。

我们想强调，指数 $m$ 是量子化的。指数的量子化与下述事实直接相关：该指数与电子的统计性质联系在一起（见式 (7.4.8)）。因此，该指数是一个拓扑数，与电子相互作用、边缘势等无关。尽管该指数是边缘态的性质，改变它的唯一途径却是通过体态中的相变。因此，这个指数可以看作一个表征\emph{体内} FQH 态拓扑序的量子数。

在动量空间中，电子传播子具有形式
\begin{equation*}
G(k,\omega)\propto\frac{(vk+\omega)^{m-1}}{\omega-vk+\mathrm{i}0^+\operatorname{sgn}(\omega)}
\end{equation*}
这个反常指数 $m$ 可以在隧穿实验中测量。由 $\operatorname{Im}G$ 可得，电子的隧穿态密度为
\begin{equation*}
N(\omega)\propto|\omega|^{m-1}
\end{equation*}
这表明，对金属--绝缘体--FQH 结，微分电导具有 $\frac{\mathrm{d}I}{\mathrm{d}V}\propto V^{m-1}$ 的形式。

\subsection*{问题 7.4.4}

证明式 (7.4.8)。（提示：先证明 $[\phi(x),\phi(y)]=\frac{\pi}{q}\operatorname{sgn}(x-y)$；译注：原文之 $q$ 疑为 $\nu$ 之误。）

\subsection*{问题 7.4.5}

\textbf{玻色化与费米化（bosonization and fermionization）}当 $m=1$ 时，本节的讨论意味着 $\nu=1$ 的 IQH 边缘态可以用自由声子理论
\begin{equation*}
H_b=2\pi v\sum_{k>0}\rho_k^\dagger\rho_k
\end{equation*}
来描述。在 7.4.1 节中，我们发现 $\nu=1$ 的 IQH 边缘态可以用自由手征费米子模型
\begin{equation*}
H_f=\sum_k vk:c_k^\dagger c_k:=\sum_{k>0}vkc_k^\dagger c_k-\sum_{k<0}vkc_kc_k^\dagger
\end{equation*}
来描述。正规排序定义为：当 $k>0$ 时 $:c_k^\dagger c_k:=c_k^\dagger c_k$；当 $k<0$ 时 $:c_k^\dagger c_k:=-c_kc_k^\dagger$。在本问题中，我们将研究这两种描述之间的直接关系。从费米子描述切换到玻色子描述称为玻色化，从玻色子描述切换到费米子描述称为费米化。
\begin{enumerate}
\item 求玻色模型前五个能级的能量及其简并度。将你的结果与问题 7.4.2 的结果作比较。
\item 令 $\rho_f(x)=:c^\dagger(x)c(x):$。在 $k$ 空间中，$\rho_{f,k}=\sum_q{}^{\prime}:c_q^\dagger c_{q+k}:$，其中求和 $\sum_q{}^{\prime}$ 限制在使 $c_{q+k}$ 与 $c_q$ 的动量都处于 $[-\Lambda,+\Lambda]$ 范围之内的那些项。正规排序后的 $\rho_f(x)$ 在基态中的平均值为零，这与玻色理论中的 $\rho$ 一致。证明 $\rho_{f,k}$ 满足与 $\rho_k$ 相同的 K--M 代数（见式 (7.4.5)，取 $\nu=1$）。（提示：对大的正 $k$ 有 $c_k^\dagger c_k=0$，对大的负 $k$ 也有 $c_k^\dagger c_k=0$。）
\item 证明 $H_f=2\pi v\sum_{k>0}\rho_{f,k}^\dagger\rho_{f,k}$。
\end{enumerate}

\subsection{边缘激发的微观理论}

\begin{keypoints}
\item K--M 代数与 Laughlin 波函数之间的关系。
\item 由等离子体类比得到的电子等时关联。
\end{keypoints}

在这一节中，我们将给出 Laughlin 态边缘激发的一个微观理论。具体起见，让我们考虑第一朗道能级中的电子气体。我们假设电子由理想哈密顿量 (7.2.5) 描述。$\nu=1/3$ 的 Laughlin 波函数
\begin{equation*}
\Psi_3(z_i)=Z^{-1/2}\prod_{i<j}(z_i-z_j)^3\prod_k\mathrm{e}^{-|z_k|^2/4l_B^2}
\tag{7.4.10}
\end{equation*}
具有零能量\footnote{这里，为方便起见，我们把能量零点移到了第零朗道能级。}，并且是理想哈密顿量的一个精确基态。（译注：原书指数中印作 $4l_B$，应为 $4l_B^2$。）在这一节中，我们将假设磁长度 $l_B=1$。式 (7.4.10) 中的 $Z$ 是归一化因子。然而，Laughlin 态 (7.4.10) 并不是唯一的零能量态。容易检验，下列类型的态都具有零能量：
\begin{equation*}
\Psi(z_i)=P(z_i)\Psi_3(z_i)
\tag{7.4.11}
\end{equation*}
其中 $P(z_i)$ 是 $z_i$ 的对称多项式。事实上，反之亦然：所有零能态都具有 (7.4.11) 的形式。这是因为，要使一个费米子态具有零能量，当任意两个电子 $i$ 与 $j$ 相互靠拢时，$\Psi$ 必须至少以 $(z_i-z_j)^3$ 的快慢消失（$(z_i-z_j)^2$ 的可能性被费米统计所排除）。由于 Laughlin 波函数只在 $z_i=z_j$ 时为零，故 $P=\Psi/\Psi_3$ 是一个有限函数。由于 $\Psi$ 与 $\Psi_3$ 都是第一朗道能级中的反对称函数，$P$ 是一个对称的全纯函数，因而只能是一个对称多项式。

在式 (7.4.11) 的所有态中，Laughlin 态描述半径最小的一个圆形液滴。所有其他的态都是 Laughlin 态液滴的形变和/或膨胀。因此，由 $P$ 生成的态对应于 Laughlin 态的边缘激发。

首先，让我们考虑零能空间（即对称多项式的空间）。我们知道，对称多项式的空间由多项式 $s_n=\sum_iz_i^n$（通过乘法与加法）生成。令 $M_0=3\frac{N(N-1)}{2}$ 为 Laughlin 态 (7.4.10) 的总角动量，则态 $\Psi$ 将具有角动量 $M=\Delta M+M_0$，其中 $\Delta M$ 是对称多项式 $P$ 的次数。由于我们只有一个零次对称多项式和一个一次对称多项式，即 $s_0=1$ 与 $s_1=\sum_iz_i$，故 $\Delta M=0,1$ 的零能态是非简并的。然而，当 $\Delta M=2$ 时，我们有两个由 $P=s_2$ 与 $P=s_1^2$ 生成的零能态。对一般的 $\Delta M$，零能态的简并度由下式给出：
\begin{equation*}
\begin{array}{cccccccc}
\Delta M: & 0 & 1 & 2 & 3 & 4 & 5 & 6\\
\text{简并度}: & 1 & 1 & 2 & 3 & 5 & 7 & 11
\end{array}
\tag{7.4.12}
\end{equation*}

这里我们想指出，式 (7.4.12) 中的简并度正是我们从宏观理论所预期的。我们知道，对于圆形液滴，角动量 $\Delta M$ 可以视为沿边缘的动量 $k=2\pi\Delta M/L$，其中 $L$ 是 FQH 液滴的周长。按照宏观理论，（中性的）边缘激发由密度算符 $\rho_k$ 生成。容易检验，由密度算符生成的边缘态对每一个 $\Delta M$ 都具有与式 (7.4.12) 相同的简并度。例如，$\Delta M=2$ 的两个态由 $\rho_{\kappa_0}^2$ 与 $\rho_{2\kappa_0}$ 生成，其中 $\kappa_0=2\pi/L$。因此，由 K--M 代数 (7.4.5) 生成的空间与对称多项式的空间是同一个空间。

现在让我们问下面这个物理问题：对称多项式是否生成所有的低能态？如果是这样，那么由上述讨论我们看到，FQH 液滴（译注：原文印作“HQ 液滴”）的所有低能激发都由 K--M 代数生成，我们就可以说式 (7.4.5) 是低能激发的一个完备理论。遗憾的是，到目前为止我们还没有上述论断的解析证明。这是因为，虽然与对称多项式所生成的态正交的态具有非零能量，但这些能量在热力学极限下是否保持有限并不清楚。能隙有可能在热力学极限下趋于零。为了解决这个问题，目前我们只能依靠数值计算。在图 7.17 中，我们给出了六个电子的系统在前二十二个轨道上、对本节开头所引入哈密顿量的能谱。零能态在 $M=45,\ldots,51$（即 $\Delta M=0,\ldots,6$）处的简并度分别求得为 $1,1,2,3,5,7,11$，与式 (7.4.12) 相符。更重要的是，我们清楚地看到，一个有限的能隙把所有其余的态与零能态分隔开。因此，数值结果意味着，Laughlin 态的所有低能边缘激发都由对称多项式或 K--M 代数 (7.4.5) 生成。

\begin{figure}[H]
\centering
\includegraphics[width=0.72\textwidth]{fig_7.17.pdf}
\caption*{图 7.17\quad 六个电子的系统在前二十二个轨道中对于哈密顿量 $H_V$ 的能谱。$M=45,\ldots,51$ 处零能态的简并度分别求得为 $1,1,2,3,5,7,11$。}
\end{figure}

在下面，我们将从 Laughlin 波函数 $\Psi_m$ 出发计算准粒子/电子的等时关联。首先，我们如下计算 $|\xi\rangle=\prod_i(1-\frac{z_i}{\xi})^n\Psi_m(\{z_i\})$ 的模：
\begin{align*}
\langle\xi|\xi\rangle&=|\xi|^{-2Nn}\frac{Z_1}{Z}\tag{7.4.13}\\
Z&=\int\prod\mathrm{d}^2z_i\ \exp\Bigl(-\beta\sum_{ij}(-m^2)\ln|z_i-z_j|-\beta\sum_k\frac{m}{4}|z_k|^2\Bigr)\\
Z_1&=\int\prod\mathrm{d}^2z_i\ \exp\Bigl[\beta\sum_{ij}m^2\ln|z_i-z_j|-\beta\sum_k\Bigl(\frac{m}{4}|z_k|^2-nm\ln|\xi-z_k|\Bigr)\Bigr]
\end{align*}
其中 $\beta=2/m$。注意，$Z$ 是由 $N$ 个“电荷”为 $m$ 的粒子组成的等离子体的配分函数，而 $Z_1$ 是该等离子体与 $\xi$ 处一个“电荷”为 $n$ 的粒子发生相互作用的配分函数。我们可以写成
\begin{equation*}
Z=\mathrm{e}^{-\beta E},\qquad Z_1=\mathrm{e}^{-\beta E_1}
\end{equation*}
其中 $E$ 与 $E_1$ 是等离子体的总能量。能量的改变 $E_1-E$ 由加在 $\xi$ 处的“电荷” $n$ 粒子与 $z_i$ 处 $N$ 个“电荷” $m$ 粒子之间的相互作用给出。当 $|\xi|$ 大时，$N$ 个“电荷” $m$ 粒子形成一个圆形液滴，我们可以把它们当作 $z=0$ 处的一个点“电荷” $mN$ 来处理。我们得到 $E_1-E=-nmN\ln|\xi|$。对于较小的 $|\xi|$，“电荷” $n$ 会改变液滴的形状，此时
\begin{equation*}
E_1-E=-nmN\ln|\xi|-\frac{n^2}{2}\Bigl[-\ln\Bigl(|\xi|-\frac{R^2}{|\xi|}\Bigr)+\ln(|\xi|)\Bigr]
\tag{7.4.14}
\end{equation*}
其中 $R=\sqrt{mN}\,l_B$ 是液滴的半径。式 (7.4.14) 中的第一项是“电荷” $n$ 与\emph{未形变}液滴之间的相互作用。第二项是液滴形变引起的修正。由于等离子体的行为像金属，这一修正可以用“电荷” $n$ 与它在 $|z|=R^2/|\xi|$ 及 $z=0$ 处的镜像之间的相互作用来表示。

在上述讨论中，我们忽略了电荷的分立性，把等离子体当作连续介质来处理。我们预期，只要 $\xi$ 不太靠近液滴，即 $|\xi|-R\gg l_B$，这一近似就能给出正确的比值 $Z_1/Z$。

由式 (7.4.14)，我们求得 $|\psi^n(\xi)\rangle$ 的模为
\begin{equation*}
\langle\xi|\xi\rangle=\Bigl(\frac{\xi\xi^*}{\xi\xi^*-R^2}\Bigr)^{n^2/m}
\tag{7.4.15}
\end{equation*}
由于内积 $\langle\tilde{\xi}|\xi\rangle$ 是 $\xi$ 的全纯函数、$\tilde{\xi}$ 的反全纯函数，式 (7.4.15) 意味着
\begin{equation*}
\langle\tilde{\xi}|\xi\rangle=\Bigl(\frac{\xi\tilde{\xi}^*}{\xi\tilde{\xi}^*-R^2}\Bigr)^{n^2/m}
\end{equation*}
注意，当 $n=m$ 时，
\begin{equation*}
\xi^{Nm}\tilde{\xi}^{*Nm}\langle\tilde{\xi}|\xi\rangle=\xi^{Nm}\tilde{\xi}^{*Nm}\Bigl(\frac{\xi\tilde{\xi}^*}{\xi\tilde{\xi}^*-R^2}\Bigr)^m
\end{equation*}
正比于 $N_e=N+1$ 电子系统边缘上等时的电子传播子 $G_e$。取 $\xi=R\mathrm{e}^{\mathrm{i}\frac{x}{R}}$ 与 $\tilde{\xi}=R$，我们得到
\begin{equation*}
G_e(x)=L^{-m}a^{m-1}\mathrm{e}^{\mathrm{i}m\left(N_e-\frac{1}{2}\right)\frac{2\pi x}{L}}\sin^{-m}(\pi x/L)
\tag{7.4.16}
\end{equation*}
当 $x$ 远小于 $L\equiv2\pi R$（译注：原文印作“$21R$”，应为液滴周长 $2\pi R$）时，上式退化为式 (7.4.9)。这里 $a$ 是 $l_B$ 量级的一个长度标度。式 (7.4.16) 可以展开如下：
\begin{align*}
G_e(x)&=L^{-m}a^{m-1}\mathrm{e}^{\mathrm{i}m(N_e-1)\frac{2\pi x}{L}}\Bigl[\sum_{n=0}^{\infty}\mathrm{e}^{-\mathrm{i}\frac{2\pi x}{L}n}\Bigr]^m\notag\\
&=L^{-m}a^{m-1}\mathrm{e}^{\mathrm{i}m(N_e-1)\frac{2\pi x}{L}}\sum_{n=0}^{\infty}C_n^{m+n-1}\,\mathrm{e}^{-\mathrm{i}\frac{2\pi x}{L}n}\\
C_n^{m+n-1}&=\frac{(n+m-1)!}{(m-1)!\,n!}\tag{7.4.17}
\end{align*}

从这个展开，我们得到角动量 $M$ 态上的电子占据数 $n_M$ 如下：
\begin{equation*}
\begin{array}{ll}
n_M=0, & \delta M=m(N_e-1)-M<0\\[4pt]
n_M=\dfrac{a^{m-1}}{L^{m-1}}C_{\delta M}^{\delta M+m-1}, & \delta M\geqslant0
\end{array}
\tag{7.4.18}
\end{equation*}
我们看到，费米边缘的准确位置处在最后一个部分占据的单粒子轨道上，即处在角动量 $m(N_e-1)$ 处（或 $k_F=\sqrt{m(N_e-1)}/l_B$ 处）。注意，当 $m(N_e-1)-M\gg m$ 时，我们有 $n_M\propto(m(N_e-1)-M)^{m-1}$。用沿边缘的动量来表示，我们有 $n_k\propto(k_F-k)^{m-1}$。我们发现，与费米液体不同，占据数 $n_k$ 在费米动量处没有跳跃。

我们想指出，式 (7.4.16) 只在 $x$ 远大于磁长度 $l_B$ 时才正确。因此，式 (7.4.18) 只在 $m(N_e-1)-M\ll\sqrt{N_e}$ 时才有效。

如果边缘激发的色散是线性的，那么在低能下 $G_e$ 只能依赖于 $x-vt$。我们立刻看到，含时的电子格林函数为
\begin{equation*}
G_e(x,t)=L^{-m}a^{m-1}\mathrm{e}^{\mathrm{i}m\left(N-\frac{1}{2}\right)\frac{2\pi(x-vt)}{L}}\sin^{-m}[\pi(x-vt)/L]
\end{equation*}
当 $m=1$ 时，上式成为一维自由费米子的电子格林函数（对费米点附近的 $k$）。

\subsection{流体动力学方法——2/5 与 2/3 态}

\begin{keypoints}
\item 分级 FQH 液体的边缘态。
\item 电子算符与准粒子算符的结构。
\item 当且仅当边缘激发能够沿相反方向传播时，边缘相互作用才能修改电子/准粒子传播子的指数。
\item 当 $\nu_1=-\nu_2=1$ 时，本节的讨论描述一维 Tomonaga--Luttinger 液体。
\end{keypoints}

在这一节中，我们将用流体动力学方法研究第二级分级态的边缘结构。我们将以 2/5 态与 2/3 态为例。特别地，我们将研究分级态边缘上电子算符与准粒子算符的结构。我们还将看到，2/3 态包含两支沿相反方向传播的边缘模式，这是相当反直觉的。

首先，让我们考虑 $\nu=\frac{2}{5}$ 的 FQH 态。按照分级理论，$\nu=\frac{2}{5}$ 的 FQH 态是在 $\nu=\frac{1}{3}$ 的 FQH 态之上凝聚准粒子而生成的。因此，2/5 态包含两种不可压缩流体组分。为了明确起见，让我们考虑一个特殊的边缘势，使得 FQH 态由两个液滴组成（见图 7.14）：一个是电子凝聚体，填充因子为 $\nu_1=\frac{1}{3}$，半径为 $r_1$；另一个是准粒子凝聚体（位于 1/3 态之上），填充因子为 $\nu_2=\frac{1}{15}$（注意 $\frac{1}{3}+\frac{1}{15}=\frac{2}{5}$），半径为 $r_2<r_1$。

当 $r_1-r_2\gg l_B$ 时，两条边缘相互独立。推广 7.4.2 节的流体动力学方法，可以证明存在两支边缘激发，其低能动力学由下式描述：
\begin{equation*}
\begin{gathered}
[\rho_{Ik},\rho_{Jk'}]=\frac{\nu_I}{2\pi}k\delta_{IJ}\delta_{k+k'}\\
H=2\pi\sum_{I,k>0}\frac{v_I}{\nu_I}\rho_{I,-k}\rho_{I,k}
\end{gathered}
\tag{7.4.19}
\end{equation*}
其中 $I=1,2$ 标记两支，$v_I$ 是边缘激发的速度。在式 (7.4.19) 中，$\rho_I$ 是一维电子密度。为了使哈密顿量下有界，我们要求 $\nu_Iv_I>0$。我们发现，$\nu=\frac{2}{5}$ FQH 态的稳定性要求两个 $v_I$ 都为正。

推广 7.4.2 节的讨论，两条边缘上的电子算符为
\begin{equation*}
\Psi_I=\mathrm{e}^{\mathrm{i}\frac{1}{\nu_I}\phi_I(x)},\qquad I=1,2
\tag{7.4.20}
\end{equation*}
其中 $\partial_x\phi_I=\frac{1}{2\pi}\rho_I$。电子传播子具有形式
\begin{equation*}
\langle T(\Psi_I(x,t)\Psi_I^\dagger(0))\rangle=\frac{\mathrm{e}^{\mathrm{i}k_Ix}}{(x-v_I t)^{-1/|\nu_I|}},\qquad I=1,2
\end{equation*}
这里 $k_I=r_I/2l_B^2$。

按照分级图像，$\nu=\frac{2}{3}$ 的 FQH 态也是由两个凝聚体形成的：一个填充因子为 1 的电子凝聚体和一个填充因子为 $-\frac{1}{3}$ 的空穴凝聚体。因此，取 $(\nu_1,\nu_2)=(1,-\frac{1}{3})$，上述讨论同样可以应用于 $\nu=\frac{2}{3}$ 的 FQH 态。同样存在两支边缘激发，但只要哈密顿量是正定的，这两支现在就具有\emph{相反}的速度。

当我们把两条边缘靠拢到 $r_1-r_2\sim l_B$ 时，两支边缘激发之间的相互作用不再可以忽略。此时，哈密顿量具有形式
\begin{equation*}
H=2\pi\sum_{k>0}V_{IJ}\rho_{I,-k}\rho_{J,k}
\tag{7.4.21}
\end{equation*}

哈密顿量 (7.4.21) 可以对角化。当 $\nu_1\nu_2>0$ 时，我们可以选择
\begin{align*}
\tilde{\rho}_{1k}&=\cos(\theta)\frac{1}{\sqrt{|\nu_1|}}\rho_{1k}+\sin(\theta)\frac{1}{\sqrt{|\nu_2|}}\rho_{2k}\\
\tilde{\rho}_{2k}&=\cos(\theta)\frac{1}{\sqrt{|\nu_2|}}\rho_{2k}-\sin(\theta)\frac{1}{\sqrt{|\nu_1|}}\rho_{1k}\\
\tan(2\theta)&=2\frac{\sqrt{|\nu_1\nu_2|}V_{12}}{|\nu_1|V_{11}-|\nu_2|V_{22}}
\tag{7.4.22}
\end{align*}
可以检验，$\tilde{\rho}$ 满足
\begin{equation*}
\begin{gathered}
[\tilde{\rho}_{Ik},\tilde{\rho}_{Jk'}]=\frac{\operatorname{sgn}(\nu_I)}{2\pi}k\delta_{IJ}\delta_{k+k'}\\
H=2\pi\sum_{I,k>0}\operatorname{sgn}(\nu_I)\tilde{v}_I\tilde{\rho}_{I,-k}\tilde{\rho}_{I,k}
\end{gathered}
\tag{7.4.23}
\end{equation*}
其中边缘激发的新速度 $\tilde{v}_I$ 由下式给出：
\begin{align*}
\operatorname{sgn}(\nu_1)\tilde{v}_1&=\frac{\cos^2(\theta)}{\cos(2\theta)}|\nu_1|V_{11}-\frac{\sin^2(\theta)}{\cos(2\theta)}|\nu_2|V_{22}\\
\operatorname{sgn}(\nu_2)\tilde{v}_2&=\frac{\cos^2(\theta)}{\cos(2\theta)}|\nu_2|V_{22}-\frac{\sin^2(\theta)}{\cos(2\theta)}|\nu_1|V_{11}
\tag{7.4.24}
\end{align*}
我们看到仍然有两支边缘激发。然而在这种情况下，具有确定速度的边缘激发是内边缘上的激发与外边缘上的激发的混合。还可以证明，只要哈密顿量 (7.4.21) 下有界，两支的速度 $\tilde{v}_I$ 就总是正的。

通过反解式 (7.4.22) 把式 (7.4.20) 中的电子算符 $\Psi_I$ 用 $\tilde{\rho}_I$ 重写之后，我们可以利用式 (7.4.24) 计算它们的传播子如下：
\begin{equation*}
\langle T(\Psi_I(x,t)\Psi_I^\dagger(0))\rangle=\mathrm{e}^{\mathrm{i}k_Ix}\frac{1}{(x-\tilde{v}_1t)^{\alpha_I}}\frac{1}{(x-\tilde{v}_2t)^{\beta_I}}
\end{equation*}
其中
\begin{equation*}
(\alpha_1,\alpha_2)=\Bigl(\frac{\cos^2\theta}{|\nu_1|},\frac{\sin^2\theta}{|\nu_2|}\Bigr),\qquad
(\beta_1,\beta_2)=\Bigl(\frac{\sin^2\theta}{|\nu_1|},\frac{\cos^2\theta}{|\nu_2|}\Bigr)
\end{equation*}

然而，当两条边缘靠近到磁长度以内时，$\Psi_I$ 就不再是边缘上最普遍的电子算符了。普遍的电子算符可以包含两边缘之间的电荷转移。对 $\nu=2/5$ 的 FQH 态，内边缘与外边缘之间隔着 $\nu=\frac{1}{3}$ 的 Laughlin 态。因此，基本的电荷转移算符由下式给出：
\begin{equation*}
\eta(x)=\mathrm{e}^{\mathrm{i}(\phi_1-\frac{\nu_1}{\nu_2}\phi_2)}=(\Psi_1\Psi_2^\dagger)^{\nu_1}
\end{equation*}
它把一份 $\nu_1e=e/3$ 的电荷从外边缘转移到内边缘。普遍的电子算符于是具有形式
\begin{equation*}
\begin{gathered}
\Psi(x)=\sum_{n=-\infty}^{+\infty}c_n\psi_n(x)\\
\psi_n(x)=\Psi_1(x)\eta^n(x)
\end{gathered}
\tag{7.4.25}
\end{equation*}

为了理解这个结果，我们注意到，无论整数 $n$ 取何值，每个算符 $\psi_n$ 总是产生一个单位的局域电荷，并且是费米性算符。因此，每个 $\psi_n$ 都是边缘电子算符的一个候选者。对一个普遍的相互作用系统，边缘上的电子算符预期是不同 $\psi_n$ 的叠加，如式 (7.4.25) 所表示。注意 $\Psi_2=\psi_{-1/\nu_1}$。$\psi_n$ 的传播子可以用与上面概述类似的方法计算，结果为
\begin{equation*}
\langle T(\psi_n(x,t)\psi_m^\dagger(0))\rangle\propto\delta_{n,m}\,\mathrm{e}^{\mathrm{i}[k_1+n\nu_1(k_2-k_1)]x}\prod_I(x-\tilde{v}_It)^{-\gamma_{In}}
\end{equation*}
其中 $\gamma_{In}$ 为
\begin{align*}
\gamma_{1n}&=\Bigl[\Bigl(n+\frac{1}{|\nu_1|}\Bigr)\sqrt{|\nu_1|}\cos\theta-\frac{n\nu_1}{\nu_2}\sqrt{|\nu_2|}\sin\theta\Bigr]^2\\
\gamma_{2n}&=\Bigl[\Bigl(n+\frac{1}{|\nu_1|}\Bigr)\sqrt{|\nu_1|}\sin\theta+\frac{n\nu_1}{\nu_2}\sqrt{|\nu_2|}\cos\theta\Bigr]^2
\tag{7.4.26}
\end{align*}

从式 (7.4.25) 与上式（译注：原文作“式 (7.4.4)”，应指上面未编号的 $\psi_n$ 传播子公式）我们看到，电子传播子在分立动量 $k=k_1+n\nu_1(k_2-k_1)$ 处具有奇异性。它们类似于相互作用一维电子系统中电子传播子的 $k_F,3k_F,\ldots$ 奇异性。

对 $\nu=\frac{2}{3}$ 的 FQH 态，我们有 $\nu_1\nu_2<0$。在这种情况下，我们需要选择
\begin{align*}
\tilde{\rho}_{1k}&=\cosh(\theta)\frac{1}{\sqrt{|\nu_1|}}\rho_{1k}+\sinh(\theta)\frac{1}{\sqrt{|\nu_2|}}\rho_{2k}\\
\tilde{\rho}_{2k}&=\cosh(\theta)\frac{1}{\sqrt{|\nu_2|}}\rho_{2k}+\sinh(\theta)\frac{1}{\sqrt{|\nu_1|}}\rho_{1k}\\
\tanh(2\theta)&=2\frac{\sqrt{|\nu_1\nu_2|}V_{12}}{|\nu_1|V_{11}+|\nu_2|V_{22}}
\tag{7.4.27}
\end{align*}
来对角化哈密顿量。可以检验，$\tilde{\rho}_I$ 同样满足 K--M 代数 (7.4.23)，只是现在
\begin{align*}
\operatorname{sgn}(\nu_1)\tilde{v}_1&=\frac{\cosh^2(\theta)}{\cosh(2\theta)}|\nu_1|V_{11}-\frac{\sinh^2(\theta)}{\cosh(2\theta)}|\nu_2|V_{22}\\
\operatorname{sgn}(\nu_2)\tilde{v}_2&=\frac{\cosh^2(\theta)}{\cosh(2\theta)}|\nu_2|V_{22}-\frac{\sinh^2(\theta)}{\cosh(2\theta)}|\nu_1|V_{11}
\tag{7.4.28}
\end{align*}
同样，只要哈密顿量 $H$ 是正定的，边缘激发的速度 $\tilde{v}_I$ 就总是具有相反的符号。电子算符仍具有 (7.4.25) 的形式，其中 $\eta=(\Psi_1\Psi_2^\dagger)^{\nu_1}$。$\psi_n$ 的传播子仍由式 (7.4.4)（译注：同前，应指上面未编号的 $\psi_n$ 传播子公式）给出，其中
\begin{align*}
\gamma_{1n}&=\Bigl[\Bigl(n+\frac{1}{|\nu_1|}\Bigr)\sqrt{|\nu_1|}\cosh\theta+\frac{n\nu_1}{\nu_2}\sqrt{|\nu_2|}\sinh\theta\Bigr]^2\\
\gamma_{2n}&=\Bigl[\Bigl(n+\frac{1}{|\nu_1|}\Bigr)\sqrt{|\nu_1|}\sinh\theta+\frac{n\nu_1}{\nu_2}\sqrt{|\nu_2|}\cosh\theta\Bigr]^2
\tag{7.4.29}
\end{align*}

在等空间点处，电子算符 $\psi_n$ 的长时间关联由总指数 $g_e^{(n)}$ 控制如下：
\begin{equation*}
\langle\psi_n^\dagger(0,t)\psi_n(0,0)\rangle\sim t^{-g_e^{(n)}},\qquad g_e^{(n)}=\sum_I\gamma_{In}
\end{equation*}

等空间的电子关联决定了边缘隧穿实验中的 $I$--$V$ 曲线（见 3.7.4 节与 4.2.2 节）。指数的最小值 $g_e\equiv\operatorname{Min}(g_e^{(n)})$ 控制着电子在两边缘之间隧穿的标度性质。例如，隧穿电导在有限温度下的标度行为为
\begin{equation*}
\sigma\propto T^{2g_e-2}
\end{equation*}
对 $\nu=2/5$ 态，这里 $g_e=3$。对 $\nu=2/3$ 态，$g_e$ 依赖于两边缘之间的相互作用。

\subsection*{问题 7.4.6}

\textbf{一维相互作用费米子系统的玻色化}考虑一个自由的一维费米子模型
\begin{equation*}
H_0=\sum_k\epsilon_k:c_k^\dagger c_k:
\end{equation*}
正规排序定义为：当 $\epsilon_k>0$ 时 $:c_k^\dagger c_k:=c_k^\dagger c_k$；当 $\epsilon_k<0$ 时 $:c_k^\dagger c_k:=-c_kc_k^\dagger$。这里当 $k=\pm k_F$ 时 $\epsilon_k=0$。
\begin{enumerate}
\item 令 $\rho_{1,k}=\sum_{q\sim k_F}^{\prime}:c_q^\dagger c_{q+k}:$，其中求和 $\sum_{q\sim k_F}^{\prime}$ 限制在使 $c_{q+k}$ 与 $c_q$ 的动量都处于范围 $[-\Lambda+k_F,+\Lambda+k_F]$ 之内的那些项。令 $\rho_{2,k}=\sum_{q\sim-k_F}^{\prime}:c_q^\dagger c_{q+k}:$，其中求和 $\sum_{q\sim-k_F}^{\prime}$ 使 $c_{q+k}$ 与 $c_q$ 处于范围 $[-\Lambda-k_F,+\Lambda-k_F]$ 之内。这里 $k<\Lambda<k_F$。我们注意到，$\rho_{1,k}$ 在 $k_F$ 附近产生右行激发，而 $\rho_{2,k}$ 在 $-k_F$ 附近产生左行激发。证明 $\rho_{I,k}$ 满足 K--M 代数 (7.4.19)，其中 $\nu_1=1$，$\nu_2=-1$。证明 $H_0=2\pi v_F\sum_{k>0}(\rho_{1,k}\rho_{1,-k}-\rho_{2,-k}\rho_{2,k})$。于是一维自由费米子系统可以被玻色化。
\item 我们可以把电子算符写成 $c(x)=\psi_1(x)+\psi_2(x)$，其中 $\psi_1(x)=\sum_{-\Lambda+k_F}^{\Lambda+k_F}\mathrm{e}^{\mathrm{i}kx}c_k$ 是 $k_F$ 附近的电子算符，而 $\psi_2(x)=\sum_{-\Lambda-k_F}^{\Lambda-k_F}\mathrm{e}^{\mathrm{i}kx}c_k$ 是 $-k_F$ 附近的电子算符。试用 $\phi_I$ 表示 $c(x)$，这里 $(2\pi)^{-1}\partial_x\phi_I(x)=\rho_I(x)$。
\item 为了验证一维相互作用费米子系统也可以被玻色化，证明相互作用项 $H_V=\sum_{q,k_1,k_2}V_q\allowbreak(c_{k_1}^\dagger c_{k_1+q})\allowbreak(c_{k_2}^\dagger c_{k_2+q})$ 在限制到两个费米点附近时变为 $2\pi\sum\allowbreak V_{IJ}\rho_{I,-k}\rho_{J,k}+$ 常数。求 $V_{IJ}$。
\item 计算相互作用模型 $H_0+H_V$ 的电子格林函数 $\langle c^\dagger(x,t)c(0)\rangle$。
\end{enumerate}

\subsection{体有效理论与边缘态}

\begin{keypoints}
\item 边缘态的结构可以由以 $K$ 矩阵表征的体拓扑序直接确定。
\end{keypoints}

在这一节中，我们将从体 FQH 态的 Chern--Simons 有效理论直接导出边缘激发的宏观理论。在这种方法中，我们不依赖于 FQH 态的任何具体构造。体拓扑序与边缘态之间的关系变得十分清楚。我必须提醒读者，本节所给出的计算是非常形式的。结果的正确性不是由形式计算来保证的，而是由与其他独立计算（例如 7.4.3 节中的那个）的比较来保证的。

为了理解有效理论与边缘态之间的关系，让我们先考虑填充因子 $\nu=1/m$ 的最简单 FQH 态，并尝试由体有效理论重新导出 7.4.2 节的结果。这样的 FQH 态由具有作用量
\begin{equation*}
S=-\frac{m}{4\pi}\int a_\mu\partial_\nu a_\lambda\varepsilon^{\mu\nu\lambda}\mathrm{d}^3x
\tag{7.4.30}
\end{equation*}
的 U(1) Chern--Simons 理论描述。

假设我们的样品具有一个边界。为简单起见，我们假设边界是 $x$ 轴，样品覆盖下半平面。

对于带边界的 FQH 液体，有效作用量 (7.4.30) 存在一个问题：由于边界的存在，它在规范变换 $a_\mu\to a_\mu+\partial_\mu f$ 下不是不变的，即 $\Delta S=-\frac{m}{4\pi}\int_{y=0}\mathrm{d}x\mathrm{d}t\,f(\partial_ta_1-\partial_1a_0)$。为了解决这个问题，我们将把规范变换限制为在边界上为零：$f(x,y=0,t)=0$。由于这一限制，$a_\mu$ 在边界上的一些自由度变成了动力学自由度。

我们知道，有效理论 (7.4.30) 只是针对无边界的体 FQH 态导出的。这里，我们将把带有受限规范变换的式 (7.4.30) 取作为带边界的 FQH 态的有效理论的定义。这样的定义肯定是自洽的。下面我们将证明，这样的定义重现了 7.4.2 节所得到的结果。

研究规范理论动力学的一种方法是选择规范条件 $a_0=0$，并把 $a_0$ 的运动方程当作一个约束。对 Chern--Simons 理论，这样的约束成为 $f_{ij}=0$。在这个约束下，我们可以把 $a_i$ 写成 $a_i=\partial_i\phi$。把它代入式 (7.4.30)，就得到边缘上一个有效的一维理论，其作用量为（Elitzur et al., 1989）
\begin{equation*}
S_{\text{edge}}=-\frac{m}{4\pi}\int\partial_t\phi\,\partial_x\phi\,\mathrm{d}x\mathrm{d}t
\tag{7.4.31}
\end{equation*}

然而，这个方法有一个问题。容易看出，与作用量 (7.4.31) 相联系的哈密顿量为零，而且式 (7.4.31) 所描述的边界激发的速度为零。因此，这个作用量不能用来描述真实 FQH 样品中任何物理的边缘激发。FQH 态中的边缘激发总是具有有限的速度。

边缘激发出现有限速度是一种边界效应。由式 (7.3.18) 定义的体有效理论不包含边缘激发速度的信息。为了从有效理论确定边缘激发的动力学，我们必须找到一种把边缘速度的信息输入进来的办法。边缘速度必须被当作体有效理论中不包含的外部参量来处理。问题在于如何把这些参量放进理论之中。

现在让我们注意到，$a_0=0$ 这个条件并不是规范固定条件的唯一选择。更普遍的规范固定条件具有形式
\begin{equation*}
a_\tau=a_0+va_x=0
\tag{7.4.32}
\end{equation*}
这里 $a_x$ 是矢量势平行于样品边界的分量，而 $v$ 是一个具有速度量纲的参量。

方便的做法是选择满足
\begin{equation*}
\tilde{x}=x-vt,\qquad \tilde{t}=t,\qquad \tilde{y}=y.
\tag{7.4.33}
\end{equation*}
的新坐标。注意，在坐标变换下规范势 $a_\mu$ 按 $\partial/\partial x^\mu$ 变换。我们发现在新坐标下，规范场的各分量
% ——— 块边界说明（起点）———
% 书页 326（p341）末句为 "We find that in the new coordinates the components of
%     the gauge field"，跨页延续到书页 327（p342）顶部 "are given by" + 式 (7.4.34)。
% 上半句译文（"我们发现，在新坐标下规范场的各分量"）由 ch7_d 收尾；
% 本块正文以 %JOIN% 起头，从"由下式给出"+式 (7.4.34) 续起，与 ch7_d 末句缝合。
% ——— 块边界说明（终点）———
% 本块为章末块：书页 334（p349）止于问题 7.4.9，其后空白，第 7 章自然收束
% （书页 335 即第 8 章章首），故本块不设 %--D-TAIL-- 区。
% ——— 覆盖范围 ———
% §7.4.5 后部（自式 (7.4.34) 起，含式 (7.4.35)–(7.4.37)、问题 7.4.7、脚注 49）；
% §7.4.6 全节（式 (7.4.38)–(7.4.48)、问题 7.4.8、脚注 50）；
% §7.4.7 全节（图 7.18、表 7.3、式 (7.4.49)、问题 7.4.9）。
% ——————————————————————————————
由下式给出：
\begin{equation*}
\tilde{a}_{\tilde{t}}=a_{t}+va_{x},\qquad \tilde{a}_{\tilde{x}}=a_{x},\qquad \tilde{a}_{\tilde{y}}=a_{y}. \tag{7.4.34}
\end{equation*}
在新坐标下，规范固定条件就化为前面讨论过的那个条件。容易看出，在式 (7.4.33) 与式 (7.4.34) 给出的变换下，Chern--Simons 作用量的形式保持不变：
\begin{equation*}
S=-\frac{m}{4\pi}\int\mathrm{d}^{3}x\,a_{\mu}\partial_{\nu}a_{\lambda}\varepsilon^{\mu\nu\lambda}
=-\frac{m}{4\pi}\int\mathrm{d}^{3}x\,\tilde{a}_{\tilde{\mu}}\partial_{\tilde{\nu}}\tilde{a}_{\tilde{\lambda}}\varepsilon^{\tilde{\mu}\tilde{\nu}\tilde{\lambda}}
\end{equation*}
重复前面的推导，我们发现边缘作用量为
\begin{equation*}
S=-\frac{m}{4\pi}\int\mathrm{d}\tilde{t}\,\mathrm{d}\tilde{x}\,\partial_{\tilde{t}}\phi\,\partial_{\tilde{x}}\phi
\end{equation*}
用原来的物理坐标表示，上述作用量具有形式
\begin{equation*}
S=-\frac{m}{4\pi}\int\mathrm{d}t\,\mathrm{d}x\,(\partial_{t}+v\partial_{x})\phi\,\partial_{x}\phi \tag{7.4.35}
\end{equation*}
这是一个手征玻色子理论（chiral boson theory）(Gross et al., 1985; Floreanini and Jackiw, 1988)。由运动方程 $\partial_{t}\phi+v\partial_{x}\phi=0$ 可以看出，式 (7.4.35) 所描述的边缘激发具有非零的速度 $v$，并且只朝一个方向运动。

为了得到式 (7.4.35) 的量子理论，我们需要把手征玻色子理论量子化。量子化可以在动量空间中进行：
\begin{equation*}
S=\frac{m}{2\pi}\int\mathrm{d}t\sum_{k>0}\left(\mathrm{i}k\dot{\phi}_{k}\phi_{-k}-vk^{2}\phi_{k}\phi_{-k}\right).
\end{equation*}
可以证明，如果我们把 $k>0$ 的 $\phi_{k}$ 视为``坐标''，那么 $\phi_{-k}$ 将正比于与之对应的动量 $\pi_{k}=\partial S/\partial\dot{\phi}_{k}$。辨认出``坐标''与``动量''之后，我们便可以得到哈密顿量以及对易关系 $[\phi_{k},\pi_{k^{\prime}}]=\mathrm{i}\delta_{kk^{\prime}}$，它们完全定义了这个量子系统。如果引入 $\rho=\frac{1}{2\pi}\partial_{x}\phi$，我们发现该量子系统由下式描述：
\begin{equation*}
[\rho_{k},\rho_{k^{\prime}}]=\frac{k\delta_{k+k^{\prime}}}{2\pi m},\qquad
H=2\pi mv\sum_{k>0}\rho_{k}^{\dagger}\rho_{k} \tag{7.4.36}
\end{equation*}

边缘激发的速度 $v$ 通过规范固定条件进入我们的理论。注意，在受限规范变换下，具有不同 $v$ 的规范固定条件 (7.4.32) 彼此不能互相变换，它们在物理上并不等价。这与我们的假设一致：规范固定条件中的 $v$ 是一个物理参量，它实际上决定了边缘激发的速度。

哈密顿量 (7.4.36) 只有当 $vm<0$ 时才有下界。理论的自洽性要求 $v$ 与 $m$ 符号相反。因此，边缘激发速度的符号（即手性（chirality））由 Chern--Simons 项前面系数的符号决定。

由于矩阵 $K$ 可以对角化，上述结果可以容易地推广到式 (7.3.18) 所描述的一般 FQH 态。所得到的有效边缘理论具有形式
\begin{equation*}
S_{\mathrm{edge}}=\frac{1}{4\pi}\int\mathrm{d}t\,\mathrm{d}x\,\bigl[K_{IJ}\partial_{t}\phi_{I}\partial_{x}\phi_{J}-V_{IJ}\partial_{x}\phi_{I}\partial_{x}\phi_{J}\bigr] \tag{7.4.37}
\end{equation*}
哈密顿量为
\begin{equation*}
H_{\mathrm{edge}}=\frac{1}{4\pi}\int\mathrm{d}t\,\mathrm{d}x\,V_{IJ}\partial_{x}\phi_{I}\partial_{x}\phi_{J}
\end{equation*}
因此，$V$ 必须是正定矩阵。利用这个结果，可以证明 $K$ 的正本征值对应于一个左移分支，而负本征值对应于一个右移分支。

$\nu=2/5$ FQH 态的有效理论由 $K=\begin{pmatrix}3 & 2\\ 2 & 3\end{pmatrix}$ 给出。由于 $K$ 具有两个正本征值，$\nu=2/5$ FQH 态的边缘激发有两个朝同一方向运动的分支。$\nu=1-\frac{1}{n}$ FQH 态则由 $K=\begin{pmatrix}1 & 0\\ 0 & -n\end{pmatrix}$ 的有效理论描述。此时 $K$ 的两个本征值符号相反，因此边缘激发的两个分支朝相反的方向运动。

\subsection*{问题 7.4.7}
证明：量子化之后，系统 (7.4.35) 由式 (7.4.36) 描述。

\subsection{带电激发与电子传播子}

\begin{keypoints}
\item 一般 FQH 边缘上的一般电子/准粒子算符。
\end{keypoints}

我们已经研究了几种简单 FQH 液体中边缘激发的动力学，发现低能边缘激发由一个自由声子理论描述。本节我们将集中讨论一般的电荷激发。特别地，我们将对最一般的（阿贝尔）FQH 态计算电子与准粒子的传播子。关键点同样是把电子或准粒子算符用声子算符 $\rho_{I}$ 写出来。一旦做到这一点，传播子便可以容易地算出，因为在低能与长波极限下声子是自由的。

我们知道，对于式 (7.3.18) 所描述的 FQH 态，其边缘态由作用量 (7.4.37) 描述。边缘激发的希尔伯特空间构成 K--M 代数的一个表示
\begin{align*}
[\rho_{Ik},\rho_{Jk^{\prime}}]&=(K^{-1})_{IJ}\frac{1}{2\pi}k\,\delta_{k+k^{\prime}}\\
k,k^{\prime}&=\text{整数}\times\frac{2\pi}{L} \tag{7.4.38}
\end{align*}
其中 $\rho_{I}=\frac{1}{2\pi}\partial_{x}\phi_{I}$ 是 FQH 态中第 $I$ 个凝聚体的边缘密度，$I,J=1,\ldots,\kappa$，而 $\kappa$ 是 $K$ 的维数。边缘上的电子密度为
\begin{equation*}
\rho_{e}=-eq_{I}\rho_{I}
\end{equation*}
边缘激发的动力学由哈密顿量
\begin{equation*}
H=2\pi\sum_{IJ}V_{IJ}\rho_{I,-k}\rho_{J,k} \tag{7.4.39}
\end{equation*}
描述，其中 $V_{IJ}$ 是正定矩阵。

我们先尝试写出边缘上的准粒子算符 $\Psi_{\pmb{l}}$，它产生一个由 $l_{I}$ 标记的准粒子。我们知道，把准粒子插到边缘上会使第 $I$ 个凝聚体的边缘密度发生一个改变 $\delta\rho_{I}$（见式 (7.3.22)）。$\delta\rho_{I}$ 满足
\begin{equation*}
\int\mathrm{d}x\,\delta\rho_{I}=(\pmb{l}^{\top}K^{-1})_{I}.
\end{equation*}
由于 $\Psi_{\pmb{l}}$ 是一个只引起密度局域改变的局域算符，我们有
\begin{equation*}
[\rho_{I}(x),\Psi_{\pmb{l}}(x^{\prime})]=l_{J}(K^{-1})_{JI}\delta(x-x^{\prime})\Psi_{\pmb{l}}(x^{\prime}) \tag{7.4.40}
\end{equation*}
利用 Kac--Moody 代数 (7.4.38)，可以证明满足式 (7.4.40) 的准粒子算符为
\begin{equation*}
\Psi_{\pmb{l}}\propto\mathrm{e}^{\mathrm{i}\phi_{I}l_{I}} \tag{7.4.41}
\end{equation*}
准粒子 $\Psi_{\pmb{l}}$ 的电荷由对易子 $[\rho_{e},\Psi_{\pmb{l}}]$ 确定，结果为
\begin{equation*}
Q_{\pmb{l}}=-e\,\pmb{q}^{\top}K^{-1}\pmb{l} \tag{7.4.42}
\end{equation*}
由式 (7.3.24) 与式 (7.4.41)，我们看到电子算符可以写成
\begin{equation*}
\Psi_{e,L}\propto\mathrm{e}^{\mathrm{i}\sum_{I}l_{I}\phi_{I}},\qquad
l_{I}=K_{I,J}L_{J},\qquad q_{I}L_{I}=1 \tag{7.4.43}
\end{equation*}
由式 (7.4.42) 可见，上述算符携带单位电荷。$\Psi_{e,L}$ 的对易可如下计算：
\begin{align*}
\Psi_{e,L}(x)\Psi_{e,L}(x^{\prime})&=(-)^{\lambda}\Psi_{e,L}(x^{\prime})\Psi_{e,L}(x)\\
\lambda&=l_{I}(K^{-1})_{IJ}l_{J}=L_{I}K_{IJ}L_{J} \tag{7.4.44}
\end{align*}
在分级基（hierarchical basis）中，$K_{II}|_{I=1}$ 为奇数，$K_{II}|_{I>1}$ 为偶数，$\pmb{q}^{\top}=(1,0,\ldots,0)$，且 $L_{1}=1$。利用这些条件，我们可以证明 $(-)^{\lambda}=-1$。式 (7.4.43) 定义的电子算符的确是费米子算符。

由于不同 $L_{I}$ 选择下的所有算符 $\Psi_{e,L}$ 都携带单位电荷并且是费米型的，每一个 $\Psi_{e,L}$ 都可能是电子算符的候选。一般地，真正的电子算符是诸 $\Psi_{e,L}$ 的下述叠加：
\begin{equation*}
\Psi_{e}=\sum_{L}C_{L}\Psi_{e,L}
\end{equation*}
这里，当我们说边缘上存在许多不同的电子算符时，我们的真正意思是：真实的物理电子算符是这些算符的一个叠加。

利用 K--M 代数 (7.4.38) 与哈密顿量 (7.4.39)，我们可以计算一个一般准粒子算符
\begin{equation*}
\Psi_{\pmb{l}}\propto\mathrm{e}^{\mathrm{i}l_{I}\phi_{I}}
\end{equation*}
（对于 $\pmb{l}$ 的合适选择，其中包括电子算符）的传播子。首先，我们注意到变换
\begin{equation*}
\rho_{I}\rightarrow U_{IJ}\tilde{\rho}_{J}
\end{equation*}
把 $(V,K)$ 变为 $(\tilde{V},\tilde{K})=(U^{\top}VU,\,U^{\top}KU)$。选取合适的 $U$，$K$ 与 $V$ 可以被同时对角化；\footnote{由于 $V$ 是正定的对称矩阵，我们可以找到一个 $U_{1}$，它把 $V$ 变为 $V_{1}=U_{1}^{\top}VU_{1}=1$，把 $K$ 变为 $K_{1}=U_{1}^{\top}KU_{1}$。注意 $K_{1}$ 是一个对称矩阵，其本征值与 $K$ 的本征值同号（尽管它们的绝对值可以不同）。然后我们做一个正交变换把 $K_{1}$ 对角化，即 $(K_{1})_{IJ}\rightarrow(K_{2})_{IJ}=\sigma_{I}|v_{I}|^{-1}\delta_{IJ}$。在这个正交变换下，$V_{1}\rightarrow V_{2}=1$ 不变且保持对角。最后，我们用一个对角矩阵把 $K_{2}$ 缩放为 $(K_{3})_{IJ}=\sigma_{I}\delta_{IJ}$。这把 $V_{2}$ 变为 $V_{3}=|v_{I}|\delta_{IJ}$。新的 $(V_{3},K_{3})$ 便导致式 (7.4.45)。}用 $\tilde{\rho}_{I}$ 表示，式 (7.4.38) 与式 (7.4.39) 变为
\begin{align*}
[\tilde{\rho}_{Ik},\tilde{\rho}_{Jk^{\prime}}]&=\sigma_{I}\delta_{IJ}\frac{1}{2\pi}k\,\delta_{k+k^{\prime}}\\
H&=2\pi\sum_{I,k>0}|v_{I}|\,\tilde{\rho}_{I,-k}\tilde{\rho}_{I,k} \tag{7.4.45}
\end{align*}
其中 $\sigma_{I}=\pm1$ 是 $K$ 本征值的符号。由 $\tilde{\rho}_{I}$ 产生的边缘激发的速度为 $v_{I}=\sigma_{I}|v_{I}|$。

用 $\tilde{\rho}_{I}$ 表示，算符 $\Psi_{\pmb{l}}$ 具有形式
\begin{equation*}
\Psi_{\pmb{l}}\propto\mathrm{e}^{\mathrm{i}\sum_{I}\tilde{l}_{I}\tilde{\phi}_{I}},\qquad
\tilde{l}_{J}=l_{I}U_{IJ} \tag{7.4.46}
\end{equation*}
由式 (7.4.45) 与式 (7.4.46)，我们看到 $\Psi_{\pmb{l}}$ 的传播子具有以下一般形式：
\begin{equation*}
\langle\Psi_{\pmb{l}}^{\dagger}(x,t)\Psi_{\pmb{l}}(0)\rangle\propto
\mathrm{e}^{\mathrm{i}l_{I}k_{I}x}\prod_{I}(x-v_{I}t+\mathrm{i}\sigma_{I}\delta)^{-\gamma_{I}},\qquad
\gamma_{I}=\tilde{l}_{I}^{2} \tag{7.4.47}
\end{equation*}
其中 $v_{I}=\sigma_{I}|v_{I}|$ 是边缘激发的速度。右移成分的总指数减去左移成分的总指数，满足求和规则\footnote{我们用到了
$\lambda_{\pmb{l}}=\sum_{I,J,I^{\prime}}l_{I^{\prime}}U_{I^{\prime}I}\sigma_{I}(U^{\top})_{I,J}l_{J}$，以及
$(K^{-1})_{IJ}=U_{II^{\prime}}\sigma_{I^{\prime}}\delta_{I^{\prime},J^{\prime}}(U^{\top})_{J^{\prime},J}$。}
\begin{equation*}
\sum_{I}\sigma_{I}\gamma_{I}\equiv\lambda_{\pmb{l}}=\pmb{l}^{\top}K^{-1}\pmb{l} \tag{7.4.48}
\end{equation*}
我们看到，差值 $\lambda_{\pmb{l}}$ 不依赖边缘相互作用 $V$，是一个拓扑量子数。由式 (7.4.44) 与式 (7.4.48)，我们看到这个差值与 $\Psi_{\pmb{l}}$ 的统计性质直接相关。如果 $\Psi_{\pmb{l}}$ 代表电子算符（或其他费米型算符），那么 $\lambda_{\pmb{l}}$ 将是一个奇整数。从式 (7.4.47) 我们还看到，算符 $\Psi_{\pmb{l}}$ 产生的激发带有接近 $\sum_{I}l_{I}k_{I}$ 的动量。

\subsection*{问题 7.4.8}
证明式 (7.4.41) 中的 $\Psi_{\pmb{l}}$ 满足式 (7.4.40)。

\subsection{手征 Luttinger 液体的唯象后果}

\begin{keypoints}
\item 三类边缘态。
\item 长程相互作用与杂质散射的效应。
\end{keypoints}

在前四节中，我们已经证明 FQH 液体边缘上的电子构成手征 Luttinger 液体。作为手征 Luttinger 液体的特征性质之一，电子与准粒子的传播子获得如下的反常指数（anomalous exponents）：
\begin{equation*}
\langle\Psi_{e}^{\dagger}(t,x=0)\Psi_{e}(0)\rangle\sim t^{-g_{e}},\qquad
\langle\Psi_{q}^{\dagger}(t,x=0)\Psi_{q}(0)\rangle\sim t^{-g_{q}}
\end{equation*}
这些反常指数可以通过 FQH 边缘之间的隧穿实验直接测量。

\begin{figure}[H]
\centering
\includegraphics[width=0.72\textwidth]{fig_7.18.pdf}
\caption*{图 7.18\quad (a) 同一 FQH 态的两条边缘之间的准粒子隧穿。(b) 不同 FQH 态的边缘之间的电子隧穿。}
\end{figure}

以下两种情形需要分开考虑：(i) 两条 FQH 液体边缘被一个绝缘体隔开；(ii) 两条边缘被一个 FQH 液体隔开。情形 (i) 中，只有电子能够在两条边缘之间隧穿（见图 7.18(b)）；情形 (ii) 中，隔开两条边缘的 FQH 液体所支持的准粒子能够隧穿（见图 7.18(a)）。情形 (i) 的隧穿算符为 $A\propto\Psi_{e1}\Psi_{e2}^{\dagger}$，其中 $\Psi_{e1}$ 与 $\Psi_{e2}$ 是两条边缘上的电子算符；情形 (ii) 的隧穿算符具有形式 $A\propto\Psi_{q1}\Psi_{q2}^{\dagger}$，其中 $\Psi_{q1}$ 与 $\Psi_{q2}$ 是两条边缘上的准粒子算符。隧穿的物理性质可以由隧穿算符的关联函数算出，而后者又可以表示为两条边缘上电子或准粒子传播子的乘积。

令 $g=g_{e}$ 对应情形 (i)，$g=g_{q}$ 对应情形 (ii)。那么在零温下，电子与准粒子算符的反常指数导致 $\langle A(t)A(0)\rangle\propto|t|^{-2g}$，从而给出一条非线性的隧穿 $I$--$V$ 曲线（见 3.7.4 节）
\begin{equation*}
I\propto V^{2g-1}
\end{equation*}
隧穿电流的噪声谱在频率 $f=QV/h$ 处还包含一个奇异性，即
\begin{equation*}
S(f)\sim\bigl|f-\frac{QV}{h}\bigr|^{2g-1} \tag{7.4.49}
\end{equation*}
其中 $V$ 是两条边缘之间的电压差，$Q$ 是隧穿电子或隧穿准粒子的电荷。在有限温度 $T$ 下，零偏压电导也具有如下幂律依赖关系：
\begin{equation*}
\sigma\propto T^{2g-2}
\end{equation*}
我们看到，反常指数可以很容易地通过隧穿实验测量。噪声谱还进一步揭示了隧穿粒子的电荷。

指数 $g_{e}$ 与 $g_{q}$ 已在 7.4.6 节中算出。为了以简单的方式总结这些结果，最好把 FQH 边缘分成以下两类：(A) 所有边缘激发朝同一方向运动；(B) 边缘激发朝两个方向传播。

对于 A 类边缘，指数 $g_{e}$ 与 $g_{q}$ 与电子及准粒子的统计性质直接相关。在这种情形下，$g_{e}$ 与 $g_{q}$ 是普适的，不依赖电子相互作用和边缘势的细节。表 7.3 列出了一些支持 A 类边缘态的 FQH 态，以及相应的指数 $g_{e,q}$ 的数值和相应粒子的电荷（注意表中只列出了 $g_{e}$ 与 $g_{q}$ 的最小值）。

\begin{center}
\small 表 7.3\quad 具有 A 类边缘态的某些 FQH 态的 $(g_{e},g_{q},Q_{q})$\\[6pt]
\begin{tabular}{cccccc}
\toprule
FQH 态 & $\nu=1/m$ & $\nu=2/5$ & $\nu=p/(pq+1)$ & $(331)_{\nu=1/2}$ & $(332)_{\nu=2/5}$ \\
\midrule
$g_{e}$ & $m$ & $3$ & $q+1$ & $3$ & $3$ \\
$g_{q}$ & $1/m$ & $2/5$ & $p/(pq+1)$ & $3/8$ & $2/5$ \\
电荷 $Q_{q}/e$ & $1/m$ & $2/5$ & $p/(pq+1)$ & $1/4$ & $2/5$ \\
\bottomrule
\end{tabular}
\end{center}

让我们讨论一下，对于填充因子为 $\nu=\frac{p}{pq+1}$（$q=$ 偶数）的分级态，上述结果是如何得到的。这些态包括 $\nu=2/5,3/7,2/9,\ldots$ 等 FQH 态。填充因子为 $\nu=\frac{p}{pq+1}$ 的分级态，在 $\pmb{q}^{\top}=(1,1,\ldots,1)$ 的基下由 $p\times p$ 矩阵 $K=1+qC$ 描述。这里 $C$ 是伪单位矩阵（pseudo-identity matrix），即 $C_{IJ}=1$，$I,J=1,\ldots,p$，而 $K^{-1}=1-\frac{q}{pq+1}C$。由于所有边缘激发都朝同一方向运动，我们有
\begin{equation*}
\langle\Psi_{\pmb{l}}^{\dagger}(x=0,t)\Psi_{\pmb{l}}(0)\rangle\propto t^{-\lambda_{\pmb{l}}}
\end{equation*}
其中 $\lambda_{\pmb{l}}$ 由式 (7.4.48) 给出。基本准粒子由 $\pmb{l}^{\top}=(1,0,\ldots,0)$ 给出，携带电荷 $\frac{1}{pq+1}$，其传播子中的指数为 $\lambda_{\pmb{l}}=1-\frac{q}{pq+1}$。指数最小的准粒子由 $\pmb{l}^{\top}=(1,\ldots,1)$ 给出，携带电荷 $\frac{p}{pq+1}$，指数为 $\lambda_{\pmb{l}}=\frac{p}{pq+1}$，它小于 $1-\frac{q}{pq+1}$（注意我们有 $q\geqslant2$ 与 $p\geqslant1$）。在低能下，正是这样的准粒子（携带电荷 $\frac{p}{pq+1}$）主导\emph{同一} FQH 流体两条边缘之间的隧穿（见 7.4.7 节）。

电子算符由 $\Psi_{e,L}=\Psi_{\pmb{l}}$（译注：原文排印为 $\psi_{\pmb{l}}$，应为 $\Psi_{\pmb{l}}$ 之误）给出，其中 $\pmb{l}$ 满足 $\sum_{I}l_{I}=pq+1$。传播子中的指数为 $\lambda_{\pmb{l}}=\sum l_{I}^{2}-q(pq+1)$。传播子中指数最小的电子算符由 $\pmb{l}^{\top}=(q,\ldots,q,q+1)$ 给出，最小指数的值为 $\lambda_{\pmb{l}}=q+1$。在低能下，正是这样的电子算符主导两个\emph{不同} FQH 流体边缘之间的隧穿。

对于 B 类边缘，$g_{e}$ 与 $g_{q}$ 不是普适的。它们的数值依赖于电子相互作用和边缘势。

\subsection*{问题 7.4.9}
利用有限电压 $V$ 下隧穿算符的关联 $\langle A(t)A(0)\rangle$，推导式 (7.4.49)。
